\newpage
\chapter{相似标准型}
\section{多项式矩阵}
\subsection{引入}
这一章
我们将要开始研究曾经提出的问题：对于向量空间上的线性变换,怎么找到一组基使得该线性变换在这组基下的表示矩阵相对简单.代数角度则是寻找一些尽可能简单的矩阵,使得所有其他矩阵均与它们相似.

第一步,要找出矩阵在相似关系下的全系不变量.第二步,构造相似标准型.

\begin{red}
\begin{remark}
    一种关系下的全系不变量是这样的量,首先,它们在这种关系下不发生改变,并且如果全系不变量相等,则可断言关系成立.

    这也就意味着,不变量不一定是全系不变量.

    对于相抵关系,找到标准型是极其容易的,但对相似关系,这个问题要复杂得多.比如特征多项式与极小多项式,二者即使加起来也不是一个相似关系的全系不变量.
\end{remark}
\end{red}
\subsection{多项式矩阵/$\lambda$-矩阵}
\begin{formal}
\begin{definition}[$\lambda$-矩阵]\label{def:多项式矩阵}
    设$\bm{A}\left(\lambda\right)=\left(a_{ij}\left(\lambda\right)\right)_{m\times n}$,其中$a_{ij}\left(\lambda\right)$是关于未定元$\lambda$的多项式,则称$\bm{A}\left(\lambda\right)$为多项式矩阵或$\lambda$-矩阵.
\end{definition}
\end{formal}
\begin{red}
    \begin{remark}
        $\lambda$-矩阵中不能出现$\lambda^{-1}$这样的非法元素.
    \end{remark}
\end{red}
\begin{formal}
\begin{proposition}[$\lambda$-矩阵性质]\label{prop:多项式矩阵性质}
    现在需要将数值矩阵的一些性质和结果推广到多项式矩阵.设$\lambda$-矩阵$\bm{A}\left(\lambda\right)=\left(a_{ij}\left(\lambda\right)\right)_{m\times n},\bm{B}\left(\lambda\right)=\left(b_{ij}\left(\lambda\right)\right)_{k\times l}$.
    
    \begin{enumerate}[label={\textup{(\arabic*)}}]
        \item $\bm{A}\left(\lambda\right)=\bm{B}\left(\lambda\right)$当且仅当$m=k,n=l$且$a_{ij}\left(\lambda\right)=b_{ij}\left(\lambda\right),\forall 1\leqslant i\leqslant m=k,1\leqslant j\leqslant n=l$
        \item 规定同阶多项式矩阵的加法（$m=k,n=l$）：\[
        \bm{A}\left(\lambda\right)\pm\bm{B}\left(\lambda\right)=\left(a_{ij}\left(\lambda\right)\pm b_{ij}\left(\lambda\right)\right)_{m\times n}
        \]
        \item 定义数乘：$\forall k\in \mathbb{K},k\bm{A}\left(\lambda\right)=\left(ka_{ij}\left(\lambda\right)\right)_{m\times n}$
        \item 定义乘法：设$\bm{B}\left(\lambda\right)=\left(b_{ij}\left(\lambda\right)\right)_{n\times p}$,则$\bm{C}\left(\lambda\right)=\bm{A}\left(\lambda\right)\bm{B}\left(\lambda\right)=\left(c_{ij}\left(\lambda\right)\right)_{m\times p}$,其中\[
        c_{ij}\left(\lambda\right)=
        \sum_{k=1}^{n}a_{ik}\left(\lambda\right)b_{kj}\left(\lambda\right)
        \]
    \end{enumerate}
\end{proposition}
\end{formal}
\begin{formal}
\begin{definition}[$\lambda$-矩阵的初等变换]
    以下为$\lambda$-矩阵的初等行（列）变换.

    \begin{enumerate}[label={\textup{(\arabic*)}}]
        \item 不同两行（列）对换
        \item 某一行（列）乘上某个非零常数
        \item 某一行（列）乘以一个多项式加到另一行（列）上
    \end{enumerate}
\end{definition}
\end{formal}
\begin{red}
\begin{remark}
    注意,第二类变换只能是非零常数,不能是一个多项式,更不能是$\lambda^{-1}$这样的非法元素.
\end{remark}
\end{red}
\begin{formal}
\begin{definition}[初等$\lambda$-阵]\label{def:初等多项式矩阵}
    对$\bm{I}_n$依次作用三类$\lambda$-初等行变换得到的矩阵称为初等$\lambda$-阵.
    \begin{enumerate}[label={\textup{(\arabic*)}}]
        \newpage
        \item \[\bm{P}_{ij}=\begin{pmatrix}
            1&&&&&&\\
            &\ddots&&&&&\\
            &&0&\cdots&1&&\\
            &&\vdots&\ddots&\vdots&&\\
            &&1&\cdots&0&&\\
            &&&&&\ddots&\\
            &&&&&&1
        \end{pmatrix}
        \]
        \item \[
        \bm{P}_i\left(c\right)=\begin{pmatrix}
            1&&&&&&\\
            &\ddots&&&&&\\
            &&c&&&&\\
            &&&\ddots&&&\\
            &&&&1&&\\
            &&&&&\ddots&\\
            &&&&&&1
        \end{pmatrix}
        \]
        \item \[
        \bm{T}_{ij}\left(f\left(\lambda\right)\right)=\begin{pmatrix}
            1&&&&&&\\
            &\ddots&&&&&\\
            &&1&&&&\\
            &&\vdots&\ddots&&&\\
            &&f\left(\lambda\right)&\cdots&1&&\\
            &&&&&\ddots&\\
            &&&&&&1
        \end{pmatrix}
        \]
    \end{enumerate}
\end{definition}
\end{formal}
\begin{red}
    \begin{remark}
        初等数字矩阵是特殊的初等$\lambda$-矩阵.
    \end{remark}
\end{red}
\begin{formal}
\begin{theorem}[初等$\lambda$-变换与初等$\lambda$-矩阵]\label{thm:初等多项式变换与初等多项式矩阵}
    $\lambda$-矩阵的初等行（列）变换等价于左（右）乘以相应的初等$\lambda$-阵.
\end{theorem}
\end{formal}
\begin{formal}
\begin{definition}[$\lambda$-矩阵的相抵]\label{def:多项式矩阵的相抵}
    若$\bm{A}\left(\lambda\right)$经过若干次初等行（列）变换可以化为$\bm{B}\left(\lambda\right)$,则称$\bm{A}\left(\lambda\right)$与$\bm{B}\left(\lambda\right)$相抵.
\end{definition}
\end{formal}
\begin{formal}
\begin{theorem}[$\lambda$-矩阵的相抵关系]\label{thm:多项式矩阵的相抵关系}
$\lambda$-矩阵的相抵关系是等价关系.
\end{theorem}
\end{formal}
\begin{formal}
\begin{definition}[$\lambda$-矩阵的逆阵]\label{def:多项式矩阵的逆阵}
    设$n$阶$\lambda$-阵$\bm{A}\left(\lambda\right),\bm{B}\left(\lambda\right)$,若\[\bm{A}\left(\lambda\right)\bm{B}\left(\lambda\right)=\bm{B}\left(\lambda\right)\bm{A}\left(\lambda\right)=\bm{I}_n\]则称$\bm{A}\left(\lambda\right)$为可逆$\lambda$-阵,$\bm{B}\left(\lambda\right)$为$\bm{A}\left(\lambda\right)$的逆阵,记作$\bm{A}\left(\lambda\right)^{-1}.$
\end{definition}
\end{formal}
\begin{red}
\begin{remark}
    数字矩阵的结论不能够随意推广到$\lambda$-矩阵,必须经过严密的推导.比如$\lambda$-矩阵$\bm{A}\left(\lambda\right)=\begin{pmatrix}
        \lambda&0\\1&0
    \end{pmatrix},\left|\bm{A}\left(\lambda\right)\right|\neq 0$但
    $\bm{A}\left(\lambda\right)$不可逆.
\end{remark}
\end{red}
\begin{formal}
\begin{theorem}[$\lambda$-阵的乘积的逆阵]\label{thm:多项式矩阵的乘积的逆阵}
        若干可逆$\lambda$-矩阵的乘积仍然可逆且
        \[
        \left(\bm{A}_1\left(\lambda\right)\cdots\bm{A}_m\left(\lambda\right)\right)^{-1}=
        \bm{A}_m\left(\lambda\right)^{-1}\cdots\bm{A}_1\left(\lambda\right)^{-1}
        \]
\end{theorem}
\end{formal}
\begin{formal}
    \begin{theorem}[初等$\lambda$-阵的逆阵]\label{thm:初等多项式矩阵的逆阵}
        上述初等$\lambda$-阵均是可逆$\lambda$-阵.
        \[
        \bm{P}_{ij}^{-1}=\bm{P}_{ij},
        \bm{P}_i\left(c\right)^{-1}=\bm{P}_i\left(c^{-1}\right),\bm{T}_{ij}\left(f\left(\lambda\right)\right)^{-1}=\bm{T}_{ij}\left(-f\left(\lambda\right)\right)
        \]
    \end{theorem}
\end{formal}
\subsection{矩阵多项式}
\begin{formal}
\begin{definition}[矩阵多项式]\label{def:矩阵多项式}
    定义
    \[
    \bm{M}\left(\lambda\right)=\bm{M}_m\lambda^m+\bm{M}_{m-1}\lambda^{m-1}+\cdots+\bm{M}_1\lambda+\bm{M}_0
    \]其中,$\bm{M}_{i}\in M_{r\times s}\left(\mathbb{K}\right)$.这样以矩阵为系数的多项式称为矩阵多项式.
\end{definition}
\end{formal}
\begin{formal}
\begin{definition}[矩阵多项式的次数]\label{def:矩阵多项式的次数}
    若$\bm{M}_m\neq \bm{O}$,定义$\deg\,\bm{M}\left(\lambda\right)=m$.约定$\deg\,\bm{O}=-\infty.$
\end{definition}
\end{formal}
\begin{formal}
\begin{definition}[矩阵多项式的相等]\label{def:矩阵多项式的相等}
    设
    \[
    \bm{M}\left(\lambda\right)=\bm{M}_m\lambda^m+\bm{M}_{m-1}\lambda^{m-1}+\cdots+\bm{M}_1\lambda+\bm{M}_0,\bm{M}_m\neq \bm{O}
    \]和
    \[
    \bm{N}\left(\lambda\right)=\bm{N}_n\lambda^n+\bm{N}_{n-1}\lambda^{n-1}+\cdots+\bm{N}_1\lambda+\bm{N}_0,\bm{N}_n\neq \bm{O}
    \]则$\bm{M}\left(\lambda\right)=\bm{N}\left(\lambda\right)$当且仅当
    $m=n,\bm{M}_i=\bm{N}_i,\forall 0\leqslant i\leqslant m=n.$
\end{definition}
\end{formal}
\begin{red}
    \begin{remark}
        多项式矩阵和矩阵多项式这两个理论体系实质是等价的.
    \end{remark}
\end{red}
\begin{formal}
\begin{definition}[其他定义]\label{def:矩阵多项式的其他定义}
    设
    \[
    \bm{M}\left(\lambda\right)=\bm{M}_m\lambda^m+\bm{M}_{m-1}\lambda^{m-1}+\cdots+\bm{M}_1\lambda+\bm{M}_0,\bm{M}_m\neq \bm{O}
    \]和
    \[
    \bm{N}\left(\lambda\right)=\bm{N}_n\lambda^n+\bm{N}_{n-1}\lambda^{n-1}+\cdots+\bm{N}_1\lambda+\bm{N}_0,\bm{N}_n\neq \bm{O}
    \]
    \begin{enumerate}[label={\textup{(\arabic*)}}]
        \item 定义加法：不妨设$m\geqslant n$则
        \begin{align*}
            &\bm{M}\left(\lambda\right)+\bm{N}\left(\lambda\right)\\
            =&
        \bm{M}_m\lambda^m+\cdots+\bm{M}_{n+1}\lambda^{n+1}+\left(\bm{M}_n+\bm{N}_n\right)\lambda^n+\cdots+\left(\bm{M}_1+\bm{N}_1\right)\lambda+\left(\bm{M}_0+\bm{N}_0\right)
        \end{align*}
        \item 定义数乘：$\forall k\in \mathbb{K},k\bm{M}\left(\lambda\right)=\left(k\bm{M}_m\right)\lambda^m+\cdots+\left(k\bm{M}_1\right)\lambda+k\bm{M}_0$
        \item 定义乘法：\[
        \bm{M}\left(\lambda\right)\cdot \bm{N}\left(\lambda\right)=\bm{P}\left(\lambda\right)=\sum\bm{P}_k\lambda^k
        \]其中$\displaystyle
        \bm{P}_k=\sum_{i+j=k}\bm{M}_i\bm{N}_j$
    \end{enumerate}
\end{definition}
\end{formal}
\begin{red}
    \begin{remark}
        由于矩阵乘法的特殊性,矩阵多项式一般不可交换顺序,同时相乘得到的新矩阵多项式也不一定能够保持次数.
    \end{remark}
\end{red}
\begin{green}
\begin{lemma}[乘积的次数]\label{lem:矩阵多项式的乘积的次数}
    $\deg\,\left(\bm{M}\left(\lambda\right)\cdot\bm{N}\left(\lambda\right)\right)\leqslant \deg\,\bm{M}\left(\lambda\right)+\deg\,\bm{N}\left(\lambda\right)$.特别地,若$\bm{M}_m$或$\bm{N}_n$可逆则取到等号.
\end{lemma}
\begin{Proof}
因为$\bm{M}\left(\lambda\right)\cdot\bm{N}\left(\lambda\right)$的首项为$\bm{M}_m\bm{N}_n\lambda^{m+n}$可能为$\bm{O}$,但当某一个可逆时,因为可逆阵乘以一个矩阵不为零矩阵于是取到等号.
\end{Proof}
\end{green}
\begin{green}
\begin{lemma}[带余除法]\label{lem:矩阵多项式的带余除法}
    设$\lambda$-阵$\bm{M}\left(\lambda\right)$和数字矩阵$\bm{B}\in M_n\left(\mathbb{K}\right)$.那么
    \[
    \bm{M}\left(\lambda\right)=\left(\lambda\bm{I}-\bm{B}\right)\bm{Q}\left(\lambda\right)+\bm{R}
    \]
    \[
    \bm{M}\left(\lambda\right)=\bm{S}\left(\lambda\right)\left(\lambda\bm{I}-\bm{B}\right)+\bm{T}
    \]
\end{lemma}
\begin{Proof}
对$\deg\,\bm{M}\left(\lambda\right)=m$进行归纳.$m=0$取$\bm{S}\left(\lambda\right)=\bm{Q}\left(\lambda\right)=\bm{O},\bm{R}=\bm{T}=\bm{M}\left(\lambda\right)$即可.

假设$\deg\,\bm{M}\left(\lambda\right)<m$时结论成立,下证$\deg\,\bm{M}\left(\lambda\right)=m$时的情况.令$\bm{Q}_1\left(\lambda\right)=\bm{M}_m\lambda^{n-1}$则
\[
\bm{M}\left(\lambda\right)-\left(\lambda\bm{I}-\bm{B}\right)\bm{Q}_1\left(\lambda\right)=\left(\bm{B}\bm{M}_m+\bm{M}_{m-1}\right)\lambda^{m-1}+\cdots+\bm{M}_0
\]故存在$\bm{Q}_2\left(\lambda\right),\bm{R}\,\textup{s.t.}$
\[
\bm{M}\left(\lambda\right)-\left(\lambda\bm{I}-\bm{B}\right)\bm{Q}_1\left(\lambda\right)=\left(\lambda\bm{I}-\bm{B}\right)\bm{Q}_2\left(\lambda\right)+\bm{R}
\]取$\bm{Q}\left(\lambda\right)=\bm{Q}_1\left(\lambda\right)+\bm{Q}_2\left(\lambda\right).$

另一方向同理,
证毕.
\end{Proof}
\end{green}
\subsection{相似与相抵}
\begin{formal}
\begin{theorem}[相似与相抵]\label{thm:相似与相抵}
    设$\bm{A},\bm{B}\in M_n\left(\mathbb{K}\right)$,则二者相似$\bm{A}\approx  \bm{B}$当且仅当它们的特征矩阵相抵$\lambda\bm{I}_n-\bm{A}\sim\lambda\bm{I}_n-\bm{B}.$
\end{theorem}
\begin{Proof}
必要性：即二者相似则特征矩阵相抵,这是显然的.设$\bm{B}=\bm{P}^{-1}\bm{AP}$,则\[
\lambda\bm{I}_n-\bm{B}=\bm{P}^{-1}\left(\lambda\bm{I}_n-\bm{A}\right)\bm{P}
\]因为可逆阵是初等数字矩阵的乘积,后者是特殊的初等$\lambda$-矩阵,于是$\lambda\bm{I}_n-\bm{B}$与$\lambda\bm{I}_n-\bm{A}$相抵.

充分性,即$\bm{M}\left(\lambda\right)\left(\lambda\bm{I}-\bm{A}\right)\bm{N}\left(\lambda\right)=\lambda\bm{I}-\bm{B}$,其中$\bm{M}\left(\lambda\right),\bm{N}\left(\lambda\right)$是可逆$\lambda$-阵.又
$\bm{M}\left(\lambda\right)\left(\lambda\bm{I}-\bm{A}\right)=\left(\lambda\bm{I}-\bm{B}\right)\bm{N}\left(\lambda\right)^{-1}$,考虑左带余除法\[
\bm{M}\left(\lambda\right)=\left(\lambda\bm{I}-\bm{B}\right)\bm{Q}\left(\lambda\right)+\bm{R}
\]即\[
\left(\lambda\bm{I}-\bm{B}\right)\left[
    \bm{Q}\left(\lambda\right)\left(\lambda\bm{I}-\bm{A}\right)-\bm{N}\left(\lambda\right)^{-1}
\right]=-\bm{R}\left(\lambda\bm{I}-\bm{A}\right)
\]断言次数$\deg\,\left[
    \bm{Q}\left(\lambda\right)\left(\lambda\bm{I}-\bm{A}\right)-\bm{N}\left(\lambda\right)^{-1}
\right]\leqslant 0$（即为数字矩阵）.考虑反证法证明,设
$\deg\,
\left[
    \bm{Q}\left(\lambda\right)\left(\lambda\bm{I}-\bm{A}\right)-\bm{N}\left(\lambda\right)^{-1}
\right]\geqslant 1
$则上式左侧次数不低于$2$,右侧次数不高于$1$,矛盾.于是$\bm{Q}\left(\lambda\right)\left(\lambda\bm{I}-\bm{A}\right)-\bm{N}\left(\lambda\right)^{-1}=-\bm{P}\in M_n\left(\mathbb{K}\right)$.于是得到
\[
\bm{P}\lambda-\bm{BP}=\bm{R}\lambda-\bm{RA}
\]这是两个矩阵多项式的相等,则$\bm{P}=\bm{R},\bm{BP}=\bm{RA}$.于是只要证明$\bm{P}=\bm{R}$是一个可逆阵即可.

因为$\bm{P}=
\bm{Q}\left(\lambda\right)\left(\lambda\bm{I}-\bm{A}\right)-\bm{N}\left(\lambda\right)^{-1}
$,于是
\[
\bm{Q}\left(\lambda\right)\left(\lambda\bm{I}-\bm{A}\right)\bm{N}\left(\lambda\right)+\bm{P}\bm{N}\left(\lambda\right)=\bm{I}_n
\]即
\[
\bm{Q}\left(\lambda\right)\bm{M}\left(\lambda\right)^{-1}\left(\lambda\bm{I}-\bm{B}\right)+\bm{P}\bm{N}\left(\lambda\right)=\bm{I}_n
\]考虑右带余除法$\bm{N}\left(\lambda\right)=\bm{S}\left(\lambda\right)\left(\lambda\bm{I}-\bm{B}\right)+\bm{T}$,于是
\[
\left[
    \bm{Q}\left(\lambda\right)\bm{M}\left(\lambda\right)^{-1}+\bm{PS}\left(\lambda\right)
\right]\left(\lambda\bm{I}-\bm{B}\right)=\bm{I}_n-\bm{PT}
\]断言$\bm{Q}\left(\lambda\right)\bm{M}\left(\lambda\right)^{-1}+\bm{PS}\left(\lambda\right)=\bm{O}$,考虑反证法证明,设其次数非负,那么上式左侧次数不低于$1$,右侧次数不高于$0$,矛盾.于是得到$\bm{PT}=\bm{I}_n\Longrightarrow \bm{P}$可逆.

综上,得证.
\end{Proof}
\end{formal}
\newpage
\section{矩阵的法式}
\subsection{法式}
\begin{green}
    \begin{lemma}[首元素]\label{lem:首元素}
        设非零$\lambda$-阵$\bm{A}\left(\lambda\right)=\left(a_{ij}\left(\lambda\right)\right)_{m\times n}$,则$\bm{A}\left(\lambda\right)$必定相抵于$\bm{B}\left(\lambda\right)=\left(b_{ij}\left(\lambda\right)\right)_{m\times n}$,其中$b_{11}\left(\lambda\right)\neq 0$且$b_{11}\left(\lambda\right)\mid b_{ij}\left(\lambda\right),\forall 1\leqslant i\leqslant m,1\leqslant j\leqslant n.$
    \end{lemma}
    \begin{Proof}
    设$k=\min\,\left\{
        \deg\,a_{ij}\left(\lambda\right)\mid
        a_{ij}\left(\lambda\right)\neq 0,1\leqslant i\leqslant m,1\leqslant j\leqslant n
    \right\}\geqslant 0$,对$k$归纳.若$k=0$则存在$a_{ij}\left(\lambda\right)$是非零常数,则通过行列对换即可.

    设非零元次数最小值小于$k$时结论成立,下面证明等于$k$时的情况.取$a_{ij}\left(\lambda\right)$使得$\deg\,a_{ij}\left(\lambda\right)=k$则通过行列对换将其换到$\left(1,1\right)$位置.不妨设$a_{11}\left(\lambda\right)$达到非零元次数最小值.

    其一,设$\exists a_{i1}\left(\lambda\right)\,\textup{s.t.}a_{11}\left(\lambda\right)\nmid a_{i1}\left(\lambda\right)$则有带余除法
    \[
    a_{i1}\left(\lambda\right)=
    a_{11}\left(\lambda\right)q\left(\lambda\right)+r\left(\lambda\right)
    \]其中$\deg\,r\left(\lambda\right)<\deg\,a_{11}\left(\lambda\right)=k$且$r\left(\lambda\right)\neq 0.$那么此时将第一行乘上$-q\left(\lambda\right)$加到第$i$行,于是得到一个具有更低的最低次数元（$r\left(\lambda\right)$）的矩阵,根据归纳假设得证.

    其二,若$\exists a_{1j}\left(\lambda\right)\,\textup{s.t.}a_{11}\left(\lambda\right)\nmid a_{1j}\left(\lambda\right)$,同理可得.

    其三,若$a_{11}\left(\lambda\right)\mid a_{i1}\left(\lambda\right),a_{11}\left(\lambda\right)\mid a_{1j}\left(\lambda\right),\forall 1\leqslant i\leqslant m,1\leqslant j\leqslant n.$此时利用初等变换将同行同列其他元素变为零.然后考虑是否有$a_{11}\left(\lambda\right)\mid a_{ij}'\left(\lambda\right),\forall 2\leqslant i\leqslant m,2\leqslant j\leqslant n$（做过变换已不等于原矩阵元素）,若是,则证毕.

    设$\exists a_{ij}'\left(\lambda\right)\,\textup{s.t.}a_{11}\left(\lambda\right)\nmid a_{ij}'\left(\lambda\right).$此时将第$i$行加到第一行,这就是上面第一种情况,于是得证.
    \end{Proof}
\end{green}
\begin{formal}
\begin{theorem}[相抵]\label{thm:相抵}
    设$n$阶$\lambda$-矩阵$\bm{A}\left(\lambda\right)$,则其相抵于
    \[
\bm{B}\left(\lambda\right)=\mathrm{diag}\,\left\{
    d_1\left(\lambda\right),d_2\left(\lambda\right),\cdots,d_r\left(\lambda\right);0,\cdots,0
\right\}
    \]其中$d_i\left(\lambda\right)$为非零首一多项式且$d_i\left(\lambda\right)\mid d_{i+1}\left(\lambda\right),\forall 1\leqslant i\leqslant r-1$
\end{theorem}
\begin{Proof}
对阶数$n$归纳,$n=1$时显然,.设小于$n$阶时成立,来证明等于$n$阶的情形.$\bm{A}\left(\lambda\right)=\bm{O}$显然,下设$\bm{A}\left(\lambda\right)\neq \bm{O}$.

由\cref{lem:首元素}知,$\bm{A}\left(\lambda\right)$必定相抵于$\bm{B}\left(\lambda\right)$,其中$b_{11}\left(\lambda\right)\neq 0$且$b_{11}\left(\lambda\right)\mid b_{ij}\left(\lambda\right),\forall 1\leqslant i,j\leqslant n.$然后利用初等变换得到
\[
\bm{B}'\left(\lambda\right)=\begin{pmatrix}
    b_{11}\left(\lambda\right)&0&\cdots&0\\
    0&b_{22}'\left(\lambda\right)&\cdots&b_{2n}'\left(\lambda\right)\\
    \vdots&\vdots&&\vdots\\
    0&b_{n2}'\left(\lambda\right)&\cdots&b_{nn}'\left(\lambda\right)
\end{pmatrix}
\]显然$b_{11}\left(\lambda\right)\mid b_{ij}'\left(\lambda\right),\forall 2\leqslant i,j\leqslant n.$将该矩阵右下记作$\bm{B}_{n-1}'\left(\lambda\right)$

由归纳假设,存在$\bm{P}\left(\lambda\right),\bm{Q}\left(\lambda\right)$作为若干个$n-1$阶初等$\lambda$-阵的乘积,使得
\[
\bm{P}\left(\lambda\right)\bm{B}_{n-1}'\left(\lambda\right)\bm{Q}\left(\lambda\right)=\mathrm{diag}\,\left\{
    d_2\left(\lambda\right),\cdots,d_r\left(\lambda\right);0,\cdots,0
\right\}
\]
其中$d_i\left(\lambda\right)$为非零首一多项式且$d_i\left(\lambda\right)\mid d_{i+1}\left(\lambda\right),\forall 2\leqslant i\leqslant r-1$.设$b_{11}\left(\lambda\right)$的首项系数为$c$,记$d_1\left(\lambda\right)=c^{-1}b_{11}\left(\lambda\right).$作
\[
\begin{pmatrix}
    c^{-1}&0\\0&\bm{I}_{n-1}
\end{pmatrix}\begin{pmatrix}
    1&0\\0&\bm{P}\left(\lambda\right)
\end{pmatrix}\bm{B}'\left(\lambda\right)\begin{pmatrix}
    1&0\\0&\bm{Q}\left(\lambda\right)
\end{pmatrix}=
\begin{pmatrix}
    d_1\left(\lambda\right)&&&&&&\\
    &d_2\left(\lambda\right)&&&&&\\
    &&\ddots&&&&\\
    &&&d_r\left(\lambda\right)&&&\\
    &&&&0&&\\
    &&&&&\ddots&\\
    &&&&&&0
\end{pmatrix}
\]因为$d_1\left(\lambda\right)$整除$\bm{B}_{n-1}'\left(\lambda\right)$的任一元素,故$d_1\left(\lambda\right)$整除$\bm{P}\left(\lambda\right)\bm{B}_{n-1}'\left(\lambda\right)\bm{Q}\left(\lambda\right)$的任一元素,于是$d_1\left(\lambda\right)\mid d_2\left(\lambda\right)$.
\end{Proof}
\end{formal}
\begin{red}
\begin{remark}
    由\cref{thm:相抵}对于任意矩阵$\bm{A}\left(\lambda\right)_{m\times n}$,必定相抵于
    \[
    \begin{pmatrix}
        d_1\left(\lambda\right)&0&\cdots&0&\bm{O}\\
        0&d_2\left(\lambda\right)&\cdots&0&\bm{O}\\
        \vdots&\vdots&\ddots&\vdots&\vdots\\
        0&0&\cdots&d_r\left(\lambda\right)&\bm{O}\\
        \bm{O}&\bm{O}&\cdots&\bm{O}&\bm{O}
    \end{pmatrix}
    \]
\end{remark}
\end{red}
\begin{red}
\begin{remark}
    \cref{thm:相抵}得到的矩阵的主对角线上非零元素个数$r$可以称为该$\lambda$-矩阵的秩,但是$\lambda$-阵的秩要比数字矩阵的秩来得弱的多,满秩并不意味着可逆.
\end{remark}
\end{red}
\begin{formal}
\begin{definition}[法式]\label{def:法式}
    \cref{thm:相抵}得到矩阵称为$\bm{A}\left(\lambda\right)$的法式或者相抵标准型.
\end{definition}
\end{formal}
\begin{red}
\begin{remark}
    \cref{def:法式}叫他相抵标准型过早,还需要证明相抵标准型不依赖于$\lambda$-矩阵初等变换的选取.
\end{remark}
\end{red}
\begin{green}
\begin{lemma}[积的行列式]\label{lem:积的行列式}
    设$n$阶$\lambda$-矩阵$\bm{A}\left(\lambda\right),\bm{B}\left(\lambda\right)$,则
    \begin{enumerate}[label={\textup{(\arabic*)}}]
        \item $\left|\bm{A}\left(\lambda\right)\cdot\bm{B}\left(\lambda\right)\right|=\left|\bm{A}\left(\lambda\right)\right|\cdot\left|\bm{B}\left(\lambda\right)\right|$
        \item 定义伴随阵：
        \[
        \bm{A}\left(\lambda\right)\bm{A}\left(\lambda\right)^*=\bm{A}\left(\lambda\right)^*\bm{A}\left(\lambda\right)=\left|\bm{A}\left(\lambda\right)\right|\bm{I}_n
        \]
    \end{enumerate}
\end{lemma}
\begin{Proof}
\begin{enumerate}[label={\textup{(\arabic*)}}]
    \item 令$f\left(\lambda\right)=
    \left|\bm{A}\left(\lambda\right)\cdot\bm{B}\left(\lambda\right)\right|-\left|\bm{A}\left(\lambda\right)\right|\cdot\left|\bm{B}\left(\lambda\right)\right|
    $是关于$\lambda$的多项式.$\forall a\in \mathbb{K}$,$f\left(a\right)=0\Longrightarrow f\left(\lambda\right)=0.$
    \item 令$\left(f_{ij}\left(\lambda\right)\right)_{n\times n}=
    \bm{A}\left(\lambda\right)\bm{A}\left(\lambda\right)^*-\left|
        \bm{A}\left(\lambda\right)
    \right|\bm{I}_n
    $其中$f_{ij}\left(\lambda\right)\in \mathbb{K}\left[\lambda\right]$.$\forall a\in \mathbb{K},
    \left(f_{ij}\left(a\right)\right)_{n\times n}=\bm{A}\left(a\right)
    \bm{A}\left(a\right)^*-\left|\bm{A}\left(a\right)\right|\bm{I}_n=\bm{O}\Longrightarrow
    f_{ij}\left(\lambda\right)=0,\forall 1\leqslant i,j\leqslant n.$\qedhere
\end{enumerate}
\end{Proof}
\end{green}
\subsection{逆阵}
\begin{formal}
\begin{theorem}[可逆的充要条件]\label{thm:可逆的充要条件}
    设$n$阶$\lambda$-阵,则以下等价：
    \begin{enumerate}[label={\textup{(\arabic*)}}]
        \item $\bm{A}\left(\lambda\right)$是可逆$\lambda$-阵
        \item $\left|\bm{A}\left(\lambda\right)\right|$是非零常数
        \item $\bm{A}\left(\lambda\right)$的相抵标准型为单位阵
        \item $\bm{A}\left(\lambda\right)$只通过行变换或列变换就能化为$\bm{I}_n$
        \item $\bm{A}\left(\lambda\right)$是若干$\lambda$-初等阵的乘积
    \end{enumerate}
\end{theorem}
\begin{Proof}
$(1)\Longrightarrow (2)$因为存在$\bm{B}\left(\lambda\right)$使得$\bm{A}\left(\lambda\right)\bm{B}\left(\lambda\right)=\bm{B}\left(\lambda\right)\bm{A}\left(\lambda\right)=\left|\bm{A}\left(\lambda\right)\right|\bm{I}_n$,取行列式立得.

$(2)\Longrightarrow (3)$由\cref{thm:相抵},存在$\bm{P}\left(\lambda\right),\bm{Q}\left(\lambda\right)$（均为初等$\lambda$-阵的乘积）使得
\[
\bm{P}\left(\lambda\right)\bm{A}\left(\lambda\right)\bm{Q}\left(\lambda\right)
=
\mathrm{diag}\,\left\{
    d_1\left(\lambda\right),d_2\left(\lambda\right),\cdots,d_r\left(\lambda\right);0,\cdots,0
\right\}
\]
其中,$d_i\left(\lambda\right)$为非零首一多项式且$d_i\left(\lambda\right)\mid d_{i+1}\left(\lambda\right),\forall 1\leqslant i\leqslant r-1.$两边取行列式得到$r=n$.并且$d_1\left(\lambda\right)d_2\left(\lambda\right)\cdots d_n\left(\lambda\right)=c\neq 0$,在由首一得到$d_1\left(\lambda\right)=d_2\left(\lambda\right)=\cdots=d_n\left(\lambda\right)=1.$于是
$\bm{A}\left(\lambda\right)$相抵于$\bm{I}_n.$

$(3)\Longrightarrow (4)$因为\cref{thm:相抵},存在$\bm{P}\left(\lambda\right),\bm{Q}\left(\lambda\right)$（均为初等$\lambda$-阵的乘积）使得
\[
\bm{P}\left(\lambda\right)\bm{A}\left(\lambda\right)\bm{Q}\left(\lambda\right)
=\bm{I}_n
\]于是$\bm{Q}\left(\lambda\right)\bm{P}\left(\lambda\right)\bm{A}\left(\lambda\right)=\bm{I}_n$于是$\bm{A}\left(\lambda\right)$只需要行变换就可以化为相抵标准型$\bm{I}_n$,同理列变换.

$(4)\Longrightarrow (5)$考虑行变换,即
\[
\bm{P}_r\left(\lambda\right)
\cdots\bm{P}_1\left(\lambda\right)\bm{A}\left(\lambda\right)=\bm{I}_n
\]
其中$\bm{P}_i\left(\lambda\right)$都是初等$\lambda$-阵.于是
\[
\bm{A}\left(\lambda\right)=
\bm{P}_1\left(\lambda\right)^{-1}
\cdots\bm{P}_r\left(\lambda\right)^{-1}
\]

$(5)\Longrightarrow (1)$根据\cref{thm:初等多项式矩阵的逆阵},初等$\lambda$-阵本身是可逆阵,他们的乘积也一定是可逆阵.

综上,证毕.
\end{Proof}
\end{formal}
\begin{green}
\begin{corollary}[逆阵的求得]\label{cor:逆阵的求得}
    在\cref{lem:积的行列式}下
    \[
    \bm{A}\left(\lambda\right)^{-1}=\frac{1}{\left|\bm{A}\left(\lambda\right)\right|}\bm{A}\left(\lambda\right)^*
    \]
\end{corollary}
\begin{Proof}
考虑到\cref{lem:积的行列式}即得.
\end{Proof}
\end{green}
\begin{formal}
\begin{theorem}[特征矩阵]\label{thm:特征矩阵}
    设$\bm{A}\in M_n\left(\mathbb{K}\right)$,则$\lambda\bm{I}_n-\bm{A}$相抵于
    \[
    \mathrm{diag}\,\left\{
        1,\cdots,1;d_1\left(\lambda\right),\cdots,d_m\left(\lambda\right)
    \right\}
    \]其中$d_i\left(\lambda\right)$是非常数的首一多项式,且$d_1\left(\lambda\right)\mid d_2\left(\lambda\right)\mid \cdots\mid d_m\left(\lambda\right).$这就是$\lambda\bm{I}_n-\bm{A}$的法式.
\end{theorem}
\begin{Proof}
因为\cref{thm:相抵},存在$\bm{P}\left(\lambda\right),\bm{Q}\left(\lambda\right)$（均为初等$\lambda$-阵的乘积）使得
\[
\bm{P}\left(\lambda\right)\left(\lambda\bm{I}_n-\bm{A}\right)\bm{Q}\left(\lambda\right)
=
\mathrm{diag}\,\left\{
    d_1\left(\lambda\right),d_2\left(\lambda\right),\cdots,d_r\left(\lambda\right);0,\cdots,0
\right\}
\]取行列式得
$r=n$且
\[
c\left|\lambda\bm{I}_n-\bm{A}\right|=d_1\left(\lambda\right)d_2\left(\lambda\right)\cdots d_n\left(\lambda\right)
\]其中$c\neq 0.$考虑到首一,立得$c=1.$最后调整顺序即可.
\end{Proof}
\end{formal}
\begin{brown}
    \begin{example}
        在上述定理中,若$\deg\,d_i\left(\lambda\right)\geqslant 1$,则$\bm{A}=c\bm{I}_n$
    \end{example}
    \begin{Proof}
    因为$\left|\lambda\bm{I}_n-\bm{A}\right|=d_1\left(\lambda\right)\cdots d_n\left(\lambda\right)$,两边取次数得到$\displaystyle 
    n=\sum_{i=1}^{n}\deg\,d_i\left(\lambda\right)\geqslant n\Longrightarrow \deg\,d_i\left(\lambda\right)=1,\forall 1\leqslant i\leqslant n$所以$d_i\left(\lambda\right)$是相互整除的首一一次多项式,于是他们只能全部相等.于是$\lambda\bm{I}_n-\bm{A}$相抵于\[
    \mathrm{diag}\,\left\{
        \lambda-c, \cdots,\lambda-c
    \right\}=\lambda\bm{I}_n-c\bm{I}_n
    \]因为\cref{thm:相似与相抵},于是$\bm{A}\approx c\bm{I}_n$,后者两边乘上$\bm{P}^{-1}$、$\bm{P}$是交换的即$\bm{A}=c\bm{I}_n$.
    \end{Proof}
\end{brown}
\begin{red}
    \begin{remark}
        于是,上述定理中的$1$,只有在$\bm{A}$是纯量阵时才会不存在.
    \end{remark}
\end{red}
\begin{brown}
\begin{example}
    计算矩阵\newpage
    \[
    \begin{pmatrix}
        0&1&-1\\
        3&-2&0\\
        -1&1&-1
    \end{pmatrix}
    \]对应的特征矩阵的相抵标准型.
\end{example}
\begin{solution}
    \begin{align*}
        \lambda\bm{I}_3-\bm{A}&=\begin{pmatrix}
            \lambda&-1&1\\
            -3&\lambda+2&0\\
            1&-1&\lambda+1
        \end{pmatrix}\\
        &\longrightarrow
\begin{pmatrix}
   1&-1&\lambda+1\\
    -3&\lambda+2&0\\
    \lambda&-1&1
\end{pmatrix}\\
&\longrightarrow
\begin{pmatrix}
    1&-1&\lambda+1\\
    0&\lambda-1&3\lambda+3\\
    0&\lambda-1&-\lambda^2-\lambda+1
\end{pmatrix}\\
&\longrightarrow
\begin{pmatrix}
    1&0&0\\
    0&\lambda-1&3\lambda+3\\
    0&\lambda-1&-\lambda^2-\lambda+1
\end{pmatrix}\\
&\longrightarrow
\begin{pmatrix}
    1&0&0\\
    0&\lambda-1&6\\
    0&\lambda-1&-\lambda^2-4\lambda+7
\end{pmatrix}\\
&\longrightarrow
\begin{pmatrix}
    1&0&0\\
    0&1&0\\
    0&0&\left(\lambda-1\right)\left(\lambda^2+4\lambda+2\right)
\end{pmatrix}
    \end{align*}
\end{solution}
\end{brown}
\newpage
\section{不变因子}
\subsection{引入}
上文\cref{def:法式}提到任意$\lambda$-矩阵$\bm{A}\left(\lambda\right)$必定相抵于
\[
\mathrm{diag}\,\left\{
    d_1\left(\lambda\right),d_2\left(\lambda\right),\cdots,d_r\left(\lambda\right);0,\cdots,0   
\right\}
\]其中
$d_i\left(\lambda\right)$是非零首一多项式且$d_i\left(\lambda\right)\mid d_{i+1}\left(\lambda\right),\forall 1\leqslant i\leqslant r-1.$
那么现在的问题是$\bm{A}\left(\lambda\right)$和$\bm{B}\left(\lambda\right)$相抵是否等价于它们的法式相同.

一方面,如果二者法式相同,根据相抵关系是一种等价关系易得.另一方面,如果二者相抵,需要证明它们的法式相等.

在此成立的基础上,我们可以得到
\begin{enumerate}[label={\textup{(\arabic*)}}]
    \item $\lambda$-矩阵在相抵关系下的全系不变量$d_1\left(\lambda\right),d_2\left(\lambda\right),\cdots,d_r\left(\lambda\right)$
    \item 相抵标准型不依赖于初等变换的选取
\end{enumerate}
\subsection{行列式因子}
\begin{formal}
\begin{definition}[行列式因子]\label{def:行列式因子}
    设$n$阶$\lambda$-阵$\bm{A}\left(\lambda\right)$,$\forall 1\leqslant k\leqslant n$,定义$\bm{A}\left(\lambda\right)$的$k$阶行列式因子$D_k\left(\lambda\right)$为
    \begin{enumerate}[label={\textup{(\arabic*)}}]
        \item 若$\bm{A}\left(\lambda\right)$的所有$k$阶子式均为零,则定义$D_k\left(\lambda\right)=0$
        \item 若$\bm{A}\left(\lambda\right)$存在$k$阶子式不为零,则定义$D_k\left(\lambda\right)$为$\bm{A}\left(\lambda\right)$的所有$k$阶子式的最大公因式（非零首一多项式）
    \end{enumerate}
\end{definition}
\end{formal}
\begin{brown}
\begin{example}
    考虑$\bm{\varLambda}=\mathrm{diag}\,\left\{d_1\left(\lambda\right),d_2\left(\lambda\right),\cdots,d_r\left(\lambda\right);0,\cdots,0\right\}$,其中$d_i\left(\lambda\right)$为非零首一多项式且$d_1\left(\lambda\right)\mid d_2\left(\lambda\right)\mid \cdots\mid d_r\left(\lambda\right).$容易得到
    \[
    D_1\left(\lambda\right)=d_1\left(\lambda\right),D_2\left(\lambda\right)=\left(d_i\left(\lambda\right)d_j\left(\lambda\right)\right)_{1\leqslant i<j\leqslant r}=d_1\left(\lambda\right)d_2\left(\lambda\right)
    \]
    \[
    D_r\left(\lambda\right)=d_1\left(\lambda\right)d_2\left(\lambda\right)\cdots d_r\left(\lambda\right),D_{r+1}\left(\lambda\right)=\cdots=D_n\left(\lambda\right)=0
    \]
\end{example}
\end{brown}
\begin{green}
\begin{lemma}[非零行列式因子的整除关系]\label{lem:非零行列式因子的整除关系}
    设$\bm{A}\left(\lambda\right)$的所有非零行列式因子为$D_1\left(\lambda\right),D_2\left(\lambda\right),\cdots,D_r\left(\lambda\right)$,则$D_1\left(\lambda\right)\mid D_2\left(\lambda\right)\mid \cdots\mid D_r\left(\lambda\right)$.
\end{lemma}
\begin{Proof}
任取$\bm{A}\left(\lambda\right)$的$i+1$阶子式$M_{i+1}$且
\[
M_{i+1}=a_{11}\left(\lambda\right)M_{11}-a_{12}\left(\lambda\right)M_{12}+\cdots+\left(-1\right)^{i+2}a_{1,i+1}\left(\lambda\right)M_{1,i+1}
\]由定义有$D_i\left(\lambda\right)\mid M_{1j},\forall 1\leqslant j\leqslant i+1\Longrightarrow
D_i\left(\lambda\right)\mid M_{i+1}$又由定义有$D_{i}\left(\lambda\right)\mid D_{i+1}\left(\lambda\right).$证毕.
\end{Proof}
\end{green}
\subsection{不变因子}
\begin{formal}
\begin{definition}[不变因子]\label{def:不变因子}
    设$\bm{A}\left(\lambda\right)$的非零行列式因子为$D_1\left(\lambda\right),D_2\left(\lambda\right),\cdots,D_r\left(\lambda\right)$,则称$\bm{A}\left(\lambda\right)$的不变因子为
    \[
    \begin{cases*}
    g_1\left(\lambda\right):=D_1\left(\lambda\right)\\
    g_2\left(\lambda\right):=\cfrac{D_2\left(\lambda\right)}{D_1\left(\lambda\right)}\\
    \cdots\cdots\cdots\cdots\\
    g_r\left(\lambda\right):=\cfrac{D_r\left(\lambda\right)}{D_{r-1}\left(\lambda\right)}
    \end{cases*}
    \]也可以一起称为不变因子组.
\end{definition}
\end{formal}
\begin{brown}
\begin{example}
    在上面的例子中,不变因子组为
    \[
    d_1\left(\lambda\right),d_2\left(\lambda\right),\cdots,d_r\left(\lambda\right)
    \]
\end{example}
\end{brown}
\begin{red}
\begin{remark}
    法式的不变因子组就是主对角线上的非零元.
\end{remark}
\end{red}
\begin{formal}
\begin{theorem}[相抵的不变因子组]\label{thm:相抵的不变因子组}
   相抵$\lambda$-阵有相同的行列式因子,进而有相同的不变因子组.
\end{theorem}
\begin{Proof}
设$\bm{B}\left(\lambda\right)=\bm{P}\left(\lambda\right)\bm{A}\left(\lambda\right)\bm{Q}\left(\lambda\right)$,其中$\bm{P}\left(\lambda\right),\bm{Q}\left(\lambda\right)$是可逆$\lambda$-阵.设$\bm{A}\left(\lambda\right),\bm{B}\left(\lambda\right)$的行列式因子为$D_k\left(\lambda\right),E_k\left(\lambda\right),1\leqslant k\leqslant n.$要证
\[
D_k\left(\lambda\right)=E_k\left(\lambda\right),\forall 1\leqslant k\leqslant n
\]

根据\cref{thm:Cauchy-Binet公式}\textup{Cauchy-Binet}公式,考虑
\begin{align*}
   & \bm{B}\left(\lambda\right)\begin{pmatrix}
        i_1&i_2&\cdots&i_k\\
        j_1&j_2&\cdots&j_k
    \end{pmatrix}=
    \bm{P}\left(\lambda\right)\bm{A}\left(\lambda\right)\bm{Q}\left(\lambda\right)\begin{pmatrix}
        i_1&i_2&\cdots&i_k\\
        j_1&j_2&\cdots&j_k
    \end{pmatrix}=\\
    &\sum_{
        \begin{smallmatrix}
            1\leqslant r_1<r_2<\cdots<r_k\leqslant n\\
            1\leqslant s_1<s_2<\cdots<s_k\leqslant n
        \end{smallmatrix}
    }\bm{P}\left(\lambda\right)\begin{pmatrix}
        i_1&i_2&\cdots&i_k\\
        r_1&r_2&\cdots&r_k
    \end{pmatrix}
    \bm{A}\left(\lambda\right)\begin{pmatrix}
        r_1&r_2&\cdots&r_k\\
        s_1&s_2&\cdots&s_k
    \end{pmatrix}
    \bm{Q}\left(\lambda\right)\begin{pmatrix}
        s_1&s_2&\cdots&s_k\\
        j_1&j_2&\cdots&j_k
    \end{pmatrix}
\end{align*}
若$D_k\left(\lambda\right)=0$显然成立.则设$D_k\left(\lambda\right)\neq 0.$因为
\[
D_k\left(\lambda\right)\mid \bm{A}\left(\lambda\right)\begin{pmatrix}
    r_1&r_2&\cdots&r_k\\
    s_1&s_2&\cdots&s_k
\end{pmatrix}
\]故
\[
D_k\left(\lambda\right)\mid \bm{B}\left(\lambda\right)\begin{pmatrix}
    i_1&i_2&\cdots&i_k\\
    j_1&j_2&\cdots&j_k
\end{pmatrix}
\]不管后者是否为零,均有$D_k\left(\lambda\right)\mid E_k\left(\lambda\right)$.

总之,因为相抵关系是等价关系,我们可以得到
\[
D_k\left(\lambda\right)=0\Longleftrightarrow
E_k\left(\lambda\right)=0,\forall 1\leqslant k\leqslant n
\]及$D_k\left(\lambda\right)\neq 0,E_k\left(\lambda\right)\neq 0$时
\[
D_k\left(\lambda\right)\mid E_k\left(\lambda\right),E_k\left(\lambda\right)\mid D_k\left(\lambda\right),\forall 1\leqslant k\leqslant n
\]即存在$0\neq c\in\mathbb{K}\,\textup{s.t.}D_k\left(\lambda\right)=cE_k\left(\lambda\right)$,由首一性即得$c=1\Longrightarrow
D_k\left(\lambda\right)=E_k\left(\lambda\right)$.

随后不变因子组证明显然.综上,证毕.
\end{Proof}
\end{formal}
\begin{formal}
\begin{theorem}[法式导出的不变因子组]\label{thm:法式导出的不变因子组}
    设$\bm{A}\left(\lambda\right)$的法式为
    \[
    \mathrm{diag}\,\left\{
        d_1\left(\lambda\right),d_2\left(\lambda\right),\cdots,d_r\left(\lambda\right);0,\cdots,0
    \right\}
    \]则$\bm{A}\left(\lambda\right)$的不变因子组为
    \[
   d_1\left(\lambda\right),d_2\left(\lambda\right),\cdots,d_r\left(\lambda\right)
    \]
\end{theorem}
\begin{Proof}
由\cref{thm:相抵的不变因子组}立得.
\end{Proof}
\end{formal}
\begin{green}
\begin{corollary}[唯一确定]\label{cor:唯一确定}
    任一$\lambda$-阵的法式与不变因子相互之间唯一确定.
\end{corollary}
\end{green}
\begin{green}
\begin{corollary}[相抵与不变因子、法式]\label{cor:相抵与不变因子、法式}
    $\bm{A}\left(\lambda\right)$与$\bm{B}\left(\lambda\right)$相抵当且仅当它们的法式相同.
\end{corollary}
\begin{Proof}
一方面,法式相同时考虑到相抵关系的传递性显然.

另一方面,设$\bm{A}\left(\lambda\right)\sim \bm{\varLambda}_1,\bm{B}\left(\lambda\right)\sim\bm{\varLambda}_2$,考虑到相抵关系的传递性,$\bm{\varLambda}_1\sim\bm{\varLambda}_2$.得证.
\end{Proof}
\end{green}
\begin{green}
\begin{corollary}[不变因子与法式的不变性]\label{cor:不变因子与法式的不变性}
    \,

    \begin{enumerate}[label={\textup{(\arabic*)}}]
        \item $\lambda$-阵在相抵关系下的全系不变量是行列式因子组或不变因子组
        \item 法式的计算不依赖于初等变换的选取
    \end{enumerate}
\end{corollary}
\begin{Proof}
\begin{enumerate}[label={\textup{(\arabic*)}}]
    \item 根据\cref{thm:相抵的不变因子组},相抵的$\lambda$-阵具有相同的行列式因子,进而有相同的不变因子组.另一方面,根据\cref{cor:唯一确定},不变因子与法式相互唯一确定且根据\cref{cor:相抵与不变因子、法式},法式相同当且仅当相抵.证毕.
    \item 不妨设$\bm{A}\left(\lambda\right)$的两个法式为$\bm{\varLambda}_1,\bm{\varLambda}_2$,只要证明$\bm{\varLambda}_1=\bm{\varLambda}_2$即可.因为$\bm{\varLambda}_1\sim\bm{\varLambda}_2$,则二者有相同的法式即自身,得证.\qedhere
\end{enumerate}
\end{Proof}
\end{green}
\begin{formal}
\begin{theorem}[数字矩阵与特征矩阵]\label{thm:数字矩阵与特征矩阵}
    设$\bm{A},\bm{B}\in M_n\left(\mathbb{K}\right)$,则$\bm{A}$与$\bm{B}$相似当且仅当它们的特征矩阵$\lambda\bm{I}_n-\bm{A},\lambda\bm{I}_n-\bm{B}$具有相同的行列式因子组或不变因子组.
\end{theorem}
\begin{Proof}
根据\cref{thm:相似与相抵},数字矩阵相似当且仅当特征矩阵相抵,由\cref{cor:不变因子与法式的不变性}得证.
\end{Proof}
\end{formal}
\begin{red}
    \begin{remark}
        数字矩阵的特征矩阵的行列式因子组或不变因子组是相似关系的全系不变量.故而将$\bm{A}\in M_n\left(\mathbb{K}\right)$的特征矩阵$\lambda\bm{I}_n-\bm{A}$的行列式因子组或不变因子组称为$\bm{A}$的行列式因子组或不变因子组.
    \end{remark}
\end{red}
\begin{formal}
\begin{theorem}[相似关系的扩张]\label{thm:相似关系的扩张}
    设$\bm{A},\bm{B}\in M_n\left(\mathbb{F}\right),\mathbb{F}\subseteq \mathbb{K}$,则$\bm{A},\bm{B}$在$\mathbb{F}$上相似当且仅当$\bm{A},\bm{B}$在$\mathbb{K}$上相似.
\end{theorem}
\begin{Proof}
因为$\bm{A},\bm{B}$在$\mathbb{F}$相似即存在$\bm{P}\in M_n\left(\mathbb{F}\right)\,\textup{s.t.}\bm{B}=\bm{P}^{-1}\bm{AP}$.另一方面,$\bm{A},\bm{B}$在$\mathbb{K}$相似即存在$\bm{Q}\in M_n\left(\mathbb{K}\right)\,\textup{s.t.}\bm{B}=\bm{Q}^{-1}\bm{AQ}$.

必要性显然.下证充分性.

设$\bm{A},\bm{B}$在$\mathbb{K}$上相似,于是特征矩阵$\lambda\bm{I}_n-\bm{A},\lambda\bm{I}_n-\bm{B}$作为$\mathbb{K}$上的$\lambda$-阵有相同的行列式因子组或不变因子组.但实际上,$\bm{A},\bm{B}\in M_n\left(\mathbb{F}\right)$,因为法式与初等变换选取无关,于是求法式过程中保证初等变换取自$\mathbb{F}\left[x\right]$即得到上述法式,于是存在$\mathbb{F}$上的可逆$\lambda$-阵$\bm{P}\left(\lambda\right),\bm{Q}\left(\lambda\right),\bm{M}\left(\lambda\right),\bm{N}\left(\lambda\right)\,\textup{s.t.}$
\[
\bm{P}\left(\lambda\right)\left(\lambda\bm{I}_n-\bm{A}\right)\bm{Q}\left(\lambda\right)=\bm{\varLambda}=\bm{M}\left(\lambda\right)\left(\lambda\bm{I}_n-\bm{B}\right)\bm{N}\left(\lambda\right)
\]于是
\[
\left(\lambda\bm{I}_n-\bm{B}\right)=
\bm{M}\left(\lambda\right)^{-1}
\bm{P}\left(\lambda\right)\left(\lambda\bm{I}_n-\bm{A}\right)\bm{Q}\left(\lambda\right)
\bm{N}\left(\lambda\right)^{-1}
\]
于是两个特征矩阵作为$\mathbb{F}$上的$\lambda$-阵相似,于是$\bm{A},\bm{B}$作为$\mathbb{F}$上的矩阵相似.证毕.
\end{Proof}
\end{formal}
\begin{green}
\begin{corollary}[不变因子扩张]\label{cor:不变因子扩张}
    行列式因子组或不变因子组在基域扩张下不改变.
\end{corollary}
\end{green}
\newpage
\section{有理标准型}
\subsection{引入}
\begin{green}
\begin{lemma}[友阵]\label{lemma:友阵}
    考虑$r$阶矩阵\[
    \bm{F}=\begin{pmatrix}
        0&1&0&\cdots&0\\
        0&0&1&\cdots&0\\
        \vdots&\vdots&\vdots&&\vdots\\
        0&0&0&\cdots&1\\
        -a_r&-a_{r-1}&-a_{r-2}&\cdots&-a_1
    \end{pmatrix}
    \]\begin{enumerate}[label={\textup{(\arabic*)}}]
        \item $\bm{F}$的行列式因子组为\[
        \overbrace{1, \ldots, 1}^{r-1 \text{个}}, f\left(\lambda\right)
        \]其中\[
        f\left(\lambda\right)=\lambda^r+a_1\lambda^{r-1}+\cdots+a_{r-1}\lambda+a_r
        \]不变因子组相同
        \item $\bm{F}$的特征多项式与极小多项式相同,均为
        \[
f\left(\lambda\right)=\lambda^r+a_1\lambda^{r-1}+\cdots+a_{r-1}\lambda+a_r
\]
    \end{enumerate}
\end{lemma}
\begin{Proof}
易得
\begin{align*}
    \lambda\bm{I}_r-\bm{F}&=\begin{pmatrix}
        \lambda&-1&0&\cdots&0\\
        0&\lambda&-1&\cdots&0\\
        \vdots&\vdots&\vdots&&\vdots\\
        0&0&0&\cdots&-1\\
        a_r&a_{r-1}&a_{r-2}&\cdots&a_1+\lambda
    \end{pmatrix}
\end{align*}
有$D_r\left(\lambda\right)=\left|
\lambda\bm{I}_r-\bm{F}
\right|=f\left(\lambda\right).$再考虑$1\leqslant r\leqslant r-1$的情况,$\lambda\bm{I}_n-\bm{F}$有一个$k$阶子式$\left(-1\right)^k\Longrightarrow D_k\left(\lambda\right)=1,\forall 1\leqslant k\leqslant r-1$.

设极小多项式$m\left(\lambda\right)$,只要证$m\left(\lambda\right)=f\left(\lambda\right)$.由\cref{thm:Cayley-Hamilton定理}\textup{Cayley-Hamilton}定理知$m\left(\lambda\right)\mid f\left(\lambda\right)\Longrightarrow
\deg\,m\left(\lambda\right)\leqslant r.$一方面,若$\deg\,m\left(\lambda\right)=r\Longrightarrow m\left(\lambda\right)=f\left(\lambda\right).$另一方面,设$\deg\,m\left(\lambda\right)<r$,设\[
m\left(\lambda\right)=c_{r-1}\lambda^{r-1}+\cdots+c_1\lambda+c_0
\]其中$c_i$不完全为零.考虑标准单位行向量,则有
\[
\bm{e}_1\bm{F}=\bm{e}_2,\bm{e}_2\bm{F}=\bm{e}_3,\cdots,\bm{e}_{r-1}\bm{F}=\bm{e}_r
\]即\[
\bm{e}_1\bm{F}^i=\bm{e}_{i+1},\forall 1\leqslant i\leqslant r-1
\]由极小多项式定义有$\bm{O}=m\left(\bm{F}\right)=
c_{r-1}\bm{F}^{r-1}+\cdots+c_1\bm{F}+c_0\bm{I}_r
$同时左乘$\bm{e}_1$
\[
\bm{O}=c_{r-1}\bm{e}_1\bm{F}^{r-1}+\cdots+c_1\bm{e}_1\bm{F}+c_0\bm{e}_1\bm{I}_r
\]利用上述迭代即得
\[
\bm{O}=c_{r-1}\bm{e}_r+\cdots+c_1\bm{e}_2+c_0\bm{e}_1
\]即$\left(
c_{r-1},\cdots,c_1,c_0
\right)=\bm{0}\Longrightarrow
c_i=0,\forall 1\leqslant i\leqslant r-1
$.这与假设矛盾.故$m\left(\lambda\right)=f\left(\lambda\right)$,证毕.
\end{Proof}
\end{green}
\begin{red}
    \begin{remark}
        \cref{lemma:友阵}中实际上是多项式相伴的友阵的转置（一些书本上也可能不取转置）.一般记作$\bm{F}\left(f\left(\lambda\right)\right)=\bm{F}.$
    \end{remark}
\end{red}
\begin{green}
\begin{lemma}[相抵的另一刻画]\label{lemma:相抵的另一刻画}
    设$\lambda$-阵$\bm{A}\left(\lambda\right)$相抵于\[
    \mathrm{diag}\,\left\{
        d_1\left(\lambda\right),d_2\left(\lambda\right),\cdots,d_n\left(\lambda\right)
    \right\}
    \]$\lambda$-阵$\bm{B}\left(\lambda\right)$相抵于\[
    \mathrm{diag}\,\left\{
        d_{i_1}\left(\lambda\right),d_{i_2}\left(\lambda\right),\cdots,d_{i_n}\left(\lambda\right)
    \right\}
    \]其中$\left(
    i_1,i_2,\cdots,i_n
    \right)$是$\left(1,2,\cdots,n\right)$的一个置换,则$\bm{A}\left(\lambda\right)\sim\bm{B}\left(\lambda\right)$.
\end{lemma}
\begin{Proof}
事实上,根据相抵关系传递,只需证明两个对角阵相抵即可.

先考虑这样简单的情况：对角阵主对角线上两个元素对换时它们是否相抵.考虑如下对角阵
\[
\mathrm{diag}\,\left\{
    d_1\left(\lambda\right),\cdots,d_i\left(\lambda\right),\cdots,d_j\left(\lambda\right),\cdots,d_n\left(\lambda\right)
\right\}
\]
要将第$i$和第$j$个对角元素对换,只需先对换第$i$行和第$j$行,再对换第$i$列和第$j$列即可.对于一般的情况,只需不断地对换相邻的两个对角元素即可.证毕.
\end{Proof}
\end{green}
\subsection{有理标准型}
\begin{formal}
\begin{theorem}[各友阵组成矩阵]\label{thm:各友阵组成矩阵}
    设$\bm{A}\in M_n\left(\mathbb{K}\right)$的不变因子组为\[
    1,1,\cdots,1;
    d_1\left(\lambda\right),d_2\left(\lambda\right),\cdots,d_k\left(\lambda\right)
    \]则$\bm{A}$相似于
    \[\bm{F}=
    \mathrm{diag}\,\left\{
\bm{F}\left(d_1\left(\lambda\right)\right),\bm{F}\left(d_2\left(\lambda\right)\right),\cdots,\bm{F}\left(d_k\left(\lambda\right)\right)
    \right\}
    \]
\end{theorem}
\begin{Proof}
因为$\lambda\bm{I}_n-\bm{A}$相抵于
\[
\mathrm{diag}\,\left\{1,1,\cdots,1;
    d_1\left(\lambda\right),d_2\left(\lambda\right),\cdots,d_k\left(\lambda\right)
\right\}
\]因为他们有相同的行列式因子组,故他们有相同的$n$阶行列式因子即行列式即
\[
\left|
\lambda\bm{I}_n-\bm{A}  
\right|=d_1\left(\lambda\right)d_2\left(\lambda\right)\cdots d_k\left(\lambda\right)
\]设$\deg\,d_i\left(\lambda\right)=n_i\geqslant 1,\forall 1\leqslant i\leqslant k$,显然$n_1+n_2+\cdots+n_k.$

因为$\lambda\bm{I}-\bm{F}\left(d_i\left(\lambda\right)\right)$的法式为\[\mathrm{diag}\,\{
    \overbrace{1, \ldots, 1}^{n_i-1 \text{个}}, d_i\left(\lambda\right)
\}\]考虑到
\[
\lambda\bm{I}-\bm{F}=\begin{pmatrix}
    \lambda\bm{I}-\bm{F}\left(d_1\left(\lambda\right)\right)&&&\\
    &\lambda\bm{I}-\bm{F}\left(d_2\left(\lambda\right)\right)&&\\
    &&\ddots&\\
    &&&\lambda\bm{I}-\bm{F}\left(d_k\left(\lambda\right)\right)
\end{pmatrix}
\]通过初等变换一定得到一个对角阵
\[
\mathrm{diag}\,\left\{
    1,\cdots,1,d_1\left(\lambda\right),1,\cdots,1,d_2\left(\lambda\right),\cdots,1,\cdots,1,d_k\left(\lambda\right)
\right\}
\]显然$1$一共有$n_1-1+\cdots+n_k-1=n_1+\cdots+n_k-k=n-k$个,根据上定理,$\lambda\bm{I}-\bm{A}\sim\lambda\bm{I}-\bm{F}$.于是数字矩阵具有相似关系即$\bm{A}\approx\bm{F}$.证毕.
\end{Proof}
\end{formal}
\begin{formal}
\begin{definition}[有理标准型]\label{def:有理标准型}
    \cref{thm:各友阵组成矩阵}中的分块对角阵
    \[\bm{F}=
    \mathrm{diag}\,\left\{
\bm{F}\left(d_1\left(\lambda\right)\right),\bm{F}\left(d_2\left(\lambda\right)\right),\cdots,\bm{F}\left(d_k\left(\lambda\right)\right)
    \right\}
    \]称为$\bm{A}$的\textup{Frobenius}标准型或有理标准型.每一个$\bm{F}\left(d_i\left(\lambda\right)\right)$称为\textup{Frobenius}块.
\end{definition}
\end{formal}
\begin{brown}
\begin{example}
    设$6$阶矩阵$\bm{A}$的不变因子组为
    \[
    1,1,1,\lambda-1,\left(\lambda-1\right)^2,\left(\lambda-1\right)^2\left(\lambda+1\right)
    \]求他的有理标准型.
\end{example}
\begin{solution}
    将非常数的不变因子展开
    \[
    \left(\lambda-1\right)^2=\lambda^2-2\lambda+1,\left(\lambda-1\right)^2\left(\lambda+1\right)=\lambda^3-\lambda^2-\lambda+1
    \]
    对应\textup{Frobenius}块为其非首项系数取负倒序,得到\[
    \bm{F}=\begin{pmatrix}
        1&&&&&\\
        &0&1&&&\\
        &-1&2&&&\\
        &&&0&1&0\\
        &&&0&0&1\\
        &&&-1&1&1
    \end{pmatrix}
    \]
\end{solution}
\end{brown}
\begin{formal}
\begin{theorem}[极小多项式]\label{thm:极小多项式}
    设$\bm{A}\in M_n\left(\mathbb{K}\right)$的不变因子组为\[
    1,\cdots,1,d_1\left(\lambda\right),\cdots,d_k\left(\lambda\right)
    \]其中非零首一多项式$d_1\left(\lambda\right)\mid d_2\left(\lambda\right)\mid \cdots\mid d_k\left(\lambda\right)$,则$\bm{A}$的极小多项式为$m\left(\lambda\right)=d_k\left(\lambda\right).$
\end{theorem}
\begin{Proof}
因为$\bm{A}$相似于其有理标准型$\bm{F}=\mathrm{diag}\,\left\{
\bm{F}\left(d_1\left(\lambda\right)\right),\bm{F}\left(d_2\left(\lambda\right)\right),\cdots,\bm{F}\left(d_k\left(\lambda\right)\right)
\right\}$且相似的矩阵具有相同的极小多项式即$m_{\bm{A}}\left(\lambda\right)=m_{\bm{F}}\left(\lambda\right)$,记$\bm{F}_i=\bm{F}\left(d_i\left(\lambda\right)\right),\forall 1\leqslant i\leqslant k$,则
\[
m_{\bm{F}}\left(\lambda\right)=\left[
    m_{\bm{F}_1}\left(\lambda\right),m_{\bm{F}_2}\left(\lambda\right),\cdots,m_{\bm{F}_k}\left(\lambda\right)
\right]
\]\cref{lemma:友阵}告诉我们
\[
m_{\bm{F}_i}\left(\lambda\right)=d_i\left(\lambda\right),\forall 1\leqslant i\leqslant k
\]于是
\[
m_{\bm{A}}\left(\lambda\right)=\left[
    d_1\left(\lambda\right),d_2\left(\lambda\right),\cdots,d_k\left(\lambda\right)
\right]=d_k\left(\lambda\right)\qedhere
\]
\end{Proof}
\end{formal}
\begin{green}
\begin{corollary}[极小多项式扩张不变]
    设$\bm{A}\in M_n\left(\mathbb{F}\right)$,其极小多项式$m_{\mathbb{F}}\left(\lambda\right)$.若$\mathbb{F}\subseteq\mathbb{K}$,将$\bm{A}$视为$\mathbb{K}$上的矩阵,其极小多项式$m_{\mathbb{K}}\left(\lambda\right)$,则$m_{\mathbb{F}}\left(\lambda\right)=m_{\mathbb{K}}\left(\lambda\right).$
\end{corollary}
\begin{Proof}
由\cref{cor:不变因子扩张}知不变因子在基域扩张下不改变及\cref{thm:极小多项式}可知.
\end{Proof}
\end{green}
\begin{brown}
\begin{example}
    考虑四阶矩阵
    \[
    \bm{A}=\begin{pmatrix}
        0&0&0&0\\
        0&0&0&0\\
        0&0&0&1\\
        0&0&0&0
    \end{pmatrix},\bm{B}=\begin{pmatrix}
        0&1&0&0\\
        0&0&0&0\\
        0&0&0&1\\
        0&0&0&0
    \end{pmatrix}
    \]
\end{example}
\begin{solution}
    它们的不变因子组为
    \begin{align*}
        &\bm{A}:1,\lambda,\lambda,\lambda^2\\
        &\bm{B}:1,1,\lambda^2,\lambda^2
    \end{align*}
    在本章开始时曾说过,即使是极小多项式与特征多项式放在一起也不能成为相似关系下的全系不变量,这里就直观地给出了这个例子：极小多项式仅仅是不变因子组的最后一项,而特征多项式是不变因子组的乘积,这两个相比于整个不变因子组来说都是不完整的,都是很弱的条件,哪怕合在一起也是弱的.
\end{solution}
\end{brown}
\begin{red}
    \begin{remark}
        我们曾说要寻找尽可能简单的矩阵来描述相似关系,但是这里的\textup{Frobenius}标准型在$d_i\left(\lambda\right)$的次数比较高时其\textup{Frobenius}块的阶数也会变得很大,这让它并不足够简单,所以我们的一个想法是对$d_i\left(\lambda\right)$做一些简化如因式分解.
    \end{remark}
\end{red}
\newpage
\section{初等因子}
\begin{formal}
\begin{definition}[准素因子]\label{def:准素因子}
    设$f\left(\lambda\right)\in\mathbb{K}\left[\lambda\right],p\left(\lambda\right)$为不可约多项式,若$\exists e\in \mathbb{Z}^+\,\textup{s.t.}$\[
    p\left(\lambda\right)^e\mid f\left(\lambda\right),p\left(\lambda\right)^{e+1}\nmid f\left(\lambda\right)
    \]则称$p\left(
        \lambda
    \right)^e$为$f\left(\lambda\right)$的准素因子.
\end{definition}
\end{formal}
\begin{brown}
\begin{example}
    由\cref{cor:标准因式分解}标准因式分解\[
     f\left(\lambda\right)=cp_1\left(\lambda\right)^{e_1}p_2\left(\lambda\right)^{e_2}\cdots p_m\left(\lambda\right)^{e_m}
    \]其中$c\neq 0,p_i\left(\lambda\right)$为互异的首一不可约多项式,$e_i\geqslant 1\left(
        \forall 1\leqslant i\leqslant m 
    \right)$.则$f\left(\lambda\right)$的准素因子为
    \[
    p_1\left(\lambda\right)^{e_1},p_2\left(\lambda\right)^{e_2},\cdots,p_m\left(\lambda\right)^{e_m}
    \]因此\cref{cor:标准因式分解}也可以叫做准素分解.
\end{example}
\end{brown}
\begin{green}
\begin{lemma}[不变因子的公共因式分解]\label{lem:不变因子的公共因式分解}
    考虑$\bm{A}\in M_n\left(\mathbb{K}\right)$的非零不变因子为\[d_1\left(\lambda\right),d_2\left(\lambda\right),\cdots,d_k\left(\lambda\right)\]
    考虑\cref{cor:公共因式分解}公共因式分解
    \[
    d_1\left(\lambda\right)=p_1\left(\lambda\right)^{e_{11}}p_2\left(\lambda\right)^{e_{12}}\cdots p_t\left(\lambda\right)^{e_{1t}}
    \]\[
    d_2\left(\lambda\right)=p_1\left(\lambda\right)^{e_{21}}p_2\left(\lambda\right)^{e_{22}}\cdots p_t\left(\lambda\right)^{e_{2t}}
    \]
    \[
    \cdots\cdots\cdots\cdots
    \]
    \[
    d_k\left(\lambda\right)=p_1\left(\lambda\right)^{e_{k1}}p_2\left(\lambda\right)^{e_{k2}}\cdots p_t\left(\lambda\right)^{e_{kt}}
    \]由不变因子的整除关系知$0\leqslant e_{1j}\leqslant e_{2j}\leqslant \cdots\leqslant e_{kj},\forall 1\leqslant j\leqslant t.$
\end{lemma}
\end{green}
\begin{formal}
\begin{definition}[初等因子]\label{def:初等因子}
    若\cref{lem:不变因子的公共因式分解}中某一个$e_{ij}>0$,则称$p_{j}\left(\lambda\right)^{e_{ij}}$为$\bm{A}$的一个初等因子.$\bm{A}$的所有初等因子构成了$\bm{A}$的初等因子组.
\end{definition}
\end{formal}
\begin{green}
\begin{corollary}[不变因子与初等因子]\label{cor:不变因子与初等因子}
     等价地,$\bm{A}$的非常数不变因子的准素因子即为$\bm{A}$的初等因子.从而$\bm{A}$的初等因子组就是$\bm{A}$的非常数不变因子的准素因子全体.
\end{corollary}
\end{green}
\begin{formal}
\begin{theorem}[初等因子确定不变因子]\label{thm:初等因子确定不变因子}
    给出$\bm{A}$的初等因子组,设其中涉及到的不可约多项式为$p_1\left(\lambda\right),p_2\left(\lambda\right),\cdots,p_t\left(\lambda\right)$,在添加一些不可约多项式的零次幂后，将同一不可约多项式的幂次降幂排列：
    \[
    \begin{matrix}
        p_1\left(\lambda\right)^{e_{k1}}&p_1\left(\lambda\right)^{e_{k-1,1}}&\cdots&p_1\left(\lambda\right)^{e_{11}}\\
        p_2\left(\lambda\right)^{e_{k2}}&p_2\left(\lambda\right)^{e_{k-1,2}}&\cdots&p_2\left(\lambda\right)^{e_{12}}\\
        \vdots&\vdots&&\vdots\\
        p_t\left(\lambda\right)^{e_{kt}}&p_t\left(\lambda\right)^{e_{k-1,t}}&\cdots&p_t\left(\lambda\right)^{e_{1t}}
    \end{matrix}
    \]其中$e_{1j}\leqslant e_{2j}\leqslant \cdots\leqslant e_{kj},\forall 1\leqslant j\leqslant t.$考虑到不变因子的整除关系,令\[
    d_k\left(\lambda\right)=p_1\left(\lambda\right)^{e_{k1}}p_2\left(\lambda\right)^{e_{k2}}\cdots p_t\left(\lambda\right)^{e_{kt}}
    \]
    \[
    d_{k-1}\left(\lambda\right)=p_1\left(\lambda\right)^{e_{k-1,1}}p_2\left(\lambda\right)^{e_{k-1,2}}\cdots p_t\left(\lambda\right)^{e_{k-1,t}}
    \]\[
    \cdots\cdots\cdots\cdots
    \]
    \[
    d_1\left(\lambda\right)=p_1\left(\lambda\right)^{e_{11}}p_2\left(\lambda\right)^{e_{12}}\cdots p_t\left(\lambda\right)^{e_{1t}}
    \]于是由初等因子组确定了不变因子组.
\end{theorem}
\end{formal}
\begin{formal}
\begin{theorem}[唯一确定可逆]\label{thm:唯一确定可逆}
    由\cref{lem:不变因子的公共因式分解},\cref{def:初等因子},\cref{cor:不变因子与初等因子}和\cref{thm:初等因子确定不变因子}所表明的这一相互确定的操作是一一对应.
\end{theorem}
\end{formal}
\begin{formal}
\begin{theorem}[全系不变量]\label{thm:全系不变量}
    初等因子组是相似关系下的全系不变量.
\end{theorem}
\end{formal}
\begin{red}
\begin{remark}
    初等因子依赖于因式分解,后者严重依赖于基域的选取，于是初等因子是依赖于基域的.
\end{remark}
\end{red}
\begin{brown}
\begin{example}
    设$\bm{A}$的不变因子为\[
    1,\cdots,1,\left(\lambda-1\right)\left(\lambda^2+1\right),\left(\lambda-1\right)^2\left(\lambda^2+1\right)\left(\lambda^2-2\right)
    \]求$\bm{A}$在$\mathbb{Q},\mathbb{R},\mathbb{C}$上的初等因子.
\end{example}
\begin{solution}
    \[
    /\mathbb{Q}:\lambda-1,\left(\lambda-1\right)^2,\lambda^2+1,\lambda^2+1,\lambda^2-2
    \]
    \[
    /\mathbb{R}:\lambda-1,\left(\lambda-1\right)^2,\lambda^2+1,\lambda^2+1,\lambda-\sqrt{2},\lambda+\sqrt{2}
    \]
    \[
    /\mathbb{C}:\lambda-1,\left(\lambda-1\right)^2,\lambda-\mathrm{i},\lambda-\mathrm{i},\lambda+\mathrm{i},\lambda+\mathrm{i},\lambda-\sqrt{2},\lambda+\sqrt{2}
    \]
\end{solution}
\end{brown}
\begin{brown}
\begin{example}
    设$\bm{A}^{10}$的初等因子组为\[
    \lambda-1,\lambda-1,\left(\lambda-1\right)^2,\left(\lambda+1\right)^2,\left(\lambda+1\right)^3,\lambda-2
    \]求他的不变因子组.
\end{example}
\begin{solution}
    容易有\[
    \begin{matrix}
        \left(\lambda-1\right)^2&\lambda-1&\lambda-1\\
        \left(\lambda+1\right)^3&\left(\lambda+1\right)^2&1\\
        \lambda-2&1&1
    \end{matrix}
    \]于是\[
    d_3\left(\lambda\right)=\left(\lambda-1\right)^2\left(\lambda+1\right)^3\left(\lambda-2\right),d_2\left(\lambda\right)=\left(\lambda-1\right)\left(\lambda+1\right)^2,d_1\left(\lambda\right)=\lambda-1
    \]于是不变因子组为
    \[
    1,\cdots,1, \lambda-1,\left(\lambda-1\right)\left(\lambda+1\right)^2,\left(\lambda-1\right)^2\left(\lambda+1\right)^3\left(\lambda-2\right)
    \]其中有$7$个$1$.
\end{solution}
\end{brown}
下面简单说明一下下一节对更精细的标准型的构造.

矩阵$\bm{A}$在$\mathbb{K}$上的初等因子为\[
p_1\left(\lambda\right)^{e_1},p_2\left(\lambda\right)^{e_2},\cdots,p_m\left(\lambda\right)^{e_m}
\]其中$p_i\left(\lambda\right)$是$\mathbb{K}$上的不可约多项式,$e_i\geqslant 1.$设其不变因子组为\[
1,\cdots,1,d_1\left(\lambda\right),\cdots,d_k\left(\lambda\right)
\]则其有理标准型\[
\bm{F}=\mathrm{diag}\,\left\{
    \bm{F}_1\left(d_1\left(\lambda\right)\right),\bm{F}_2\left(d_2\left(\lambda\right)\right),\cdots,\bm{F}_k\left(d_k\left(\lambda\right)\right)
\right\}
\]而构造一个相对简单的矩阵$\bm{J}\left(p_i\left(\lambda\right)^{e_i}\right)$使得其初等因子仅有$p_i\left(\lambda\right)^{e_i}$.于是得到
\[
\bm{J}=\mathrm{diag}\,\left\{
    \bm{J}_1\left(p_1\left(\lambda\right)^{e_1}\right),\bm{J}_2\left(p_2\left(\lambda\right)^{e_2}\right),\cdots,\bm{J}_m\left(p_m\left(\lambda\right)^{e_m}\right)  
\right\}
\]然后验证其相似于$\bm{A}$，这就是下一节需要讨论的问题.
\begin{red}
\begin{remark}
    上述问题在一般数域上是超出本门课的内容的，于是我们只会讨论在复数域$\mathbb{C}$上的情况.
\end{remark}
\end{red}
\newpage
\section{Jordan标准型}
\subsection{Jordan标准型}
\begin{green}
    \begin{lemma}[Jordan块]\label{lemma:Jordan块}
        断言矩阵
        \[
        \bm{J}=\bm{J}_r\left(\lambda_0\right)=\begin{pmatrix}
           \lambda_0&1&&&\\
              &\lambda_0&1&&\\
                &&\ddots&\ddots&\\
                &&&\ddots&1\\
                &&&&\lambda_0
        \end{pmatrix}_{r\times r}
        \]有且仅有初等因子
    \[
    \left(
        \lambda-\lambda_0
    \right)^r
    \]
    \end{lemma}
    \begin{Proof}
        特征矩阵\[
        \lambda\bm{I}-\bm{J}=\begin{pmatrix}
            \lambda-\lambda_0&-1&&&\\
            &\lambda-\lambda_0&-1&&\\
            &&\ddots&\ddots&\\
            &&&\ddots&-1\\
            &&&&\lambda-\lambda_0
        \end{pmatrix}
        \]容易有\begin{align*}
            D_r\left(\lambda\right)=\left|\lambda\bm{I}-\bm{J}\right|=\left(\lambda-\lambda_0\right)^r
        \end{align*}并且$\forall 1\leqslant k<r$,存在$k$阶子式的值为$\left(-1\right)^k\Longrightarrow D_k\left(\lambda\right)=1,\forall 1\leqslant k<r.$于是行列式因子组为\[
        1,\cdots,1,\left(\lambda-\lambda_0\right)^r
        \]不变因子组相同，初等因子组为\[
        \left(\lambda-\lambda_0\right)^r\qedhere
        \]
    \end{Proof}
\end{green}
下面给出不同于法式-行列式因子-不变因子-初等因子这一过程的求得初等因子的办法.
\begin{green}
\begin{lemma}[初等因子求得]\label{lemma:初等因子求得}
    设$\lambda\bm{I}_n-\bm{A}$相抵于对角$\lambda$-阵\[
    \mathrm{diag}\,\left\{
        f_1\left(\lambda\right),f_2\left(\lambda\right),\cdots,f_n\left(\lambda\right)
    \right\}
    \]其中$f_i\left(\lambda\right)$为首一多项式.则$f_1\left(\lambda\right),f_2\left(\lambda\right),\cdots,f_n\left(\lambda\right)$的准素因子全体即为$\bm{A}$的初等因子组.
\end{lemma}
\begin{Proof}
第一步，$\forall 1\leqslant i<j\leqslant n,\left(
    f_i\left(\lambda\right),f_j\left(\lambda\right)
\right)=g\left(\lambda\right),\left[
    f_i\left(\lambda\right),f_j\left(\lambda\right)
\right]=h\left(\lambda\right).$下面证明上述对角阵相抵于\[
\mathrm{diag}\,\left\{
    f_1\left(\lambda\right),\cdots,g\left(\lambda\right),\cdots,h\left(\lambda\right),\cdots,f_n\left(\lambda\right)
\right\}
\]并且二者的主对角元的准素因子全体相同.

不妨设$\displaystyle i=1,j=2,f_1\left(\lambda\right)=g\left(\lambda\right)q\left(\lambda\right),h\left(\lambda\right)=\frac{
    f_1\left(\lambda\right)f_2\left(\lambda\right)
}{g\left(\lambda\right)}=f_2\left(\lambda\right)q\left(\lambda\right)$且存在$u\left(\lambda\right),v\left(\lambda\right)\,\textup{s.t.}$\[
g\left(\lambda\right)=f_1\left(\lambda\right)u\left(\lambda\right)+f_2\left(\lambda\right)v\left(\lambda\right)
\]考虑到\begin{align*}
    \begin{pmatrix}
        f_1\left(\lambda\right)&\\
        &f_2\left(\lambda\right)
    \end{pmatrix}&\longrightarrow
    \begin{pmatrix}
        f_1\left(\lambda\right)&g\left(\lambda\right)\\
        &f_2\left(\lambda\right)
    \end{pmatrix}\\
    &\longrightarrow
    \begin{pmatrix}
        &g\left(\lambda\right)\\
        -h\left(\lambda\right)&f_2\left(\lambda\right)
    \end{pmatrix}\\
    &\longrightarrow\begin{pmatrix}
        &g\left(\lambda\right)\\
        -h\left(\lambda\right)&
    \end{pmatrix}\\
    &\longrightarrow\begin{pmatrix}
        g\left(\lambda\right)&\\
        &h\left(\lambda\right)
    \end{pmatrix}
\end{align*}

接下来证明$f_1\left(\lambda\right),f_2\left(\lambda\right)$与$g\left(\lambda\right),h\left(\lambda\right)$的准素因子全体相同.在数域$\mathbb{K}$上做公共因式分解\begin{align*}
&f_1\left(\lambda\right)=p_1\left(\lambda\right)^{r_1}p_2\left(\lambda\right)^{r_2}\cdots p_t\left(\lambda\right)^{r_t}\\
&f_2\left(\lambda\right)=p_1\left(\lambda\right)^{s_1}p_2\left(\lambda\right)^{s_2}\cdots p_t\left(\lambda\right)^{s_t}
\end{align*}其中$p_i\left(\lambda\right)$为互异的首一不可约多项式,$r_i,s_i\geqslant 0,\forall 1\leqslant i\leqslant t.$令$e_i=\min\,\left\{r_i,s_i\right\},k_i=\max\,\left\{r_i,s_i\right\}.$于是
\begin{align*}
    &g\left(\lambda\right)=p_1\left(\lambda\right)^{e_1}p_2\left(\lambda\right)^{e_2}\cdots p_t\left(\lambda\right)^{e_t}\\
    &h\left(\lambda\right)=p_1\left(\lambda\right)^{k_1}p_2\left(\lambda\right)^{k_2}\cdots p_t\left(\lambda\right)^{k_t}
\end{align*}于是$f_1\left(\lambda\right),f_2\left(\lambda\right)$与$g\left(\lambda\right),h\left(\lambda\right)$的准素因子全体相同.

第二步,反复利用第一步,将上述对角阵变为法式.依次将$\left(1,1\right)$元与第$\left(2,2\right)$元,第$\left(3,3\right)$元,$\cdots$,第$\left(n,n\right)$元进行第一步的操作,从而第$\left(1,1\right)$元的因式分解中不同多项式的幂次均为最小.

一般地,设$\left(1,1\right),\left(2,2\right),\cdots,\left(i-1,i-1\right)$元已经做好了,那么让第$\left(i,i\right)$元依次与$\left(i+1,i+1\right),\left(i+2,i+2\right),\cdots,\left(n,n\right)$进行第一步的操作.最终得到对角阵\[
\mathrm{diag}\,\left\{
    d_1\left(\lambda\right),d_2\left(\lambda\right),\cdots,d_n\left(\lambda\right)
\right\}
\]由第一步保证$d_1\left(\lambda\right)\mid d_2\left(\lambda\right)\mid \cdots\mid d_n\left(\lambda\right)$,根据\cref{def:法式},这就是法式.

于是该对角阵与法式拥有相同的准素因子组即相同的初等因子组.
\end{Proof}
\end{green}
\begin{brown}
\begin{example}
    设$\lambda\bm{I}-\bm{A}$相抵于\[
    \begin{pmatrix}
        1&&&&\\
        &\left(\lambda-1\right)^2\left(\lambda+2\right)&&&\\
        &&\lambda+2&&\\
        &&&1&\\
        &&&&\lambda-1
    \end{pmatrix}
    \]根据\cref{lemma:初等因子求得}可知初等因子组为\[
    \lambda-1,\left(\lambda-1\right)^2,\lambda+2,\lambda+2
    \]
\end{example}
\end{brown}
\begin{green}
\begin{lemma}[Jordan块的初等因子]\label{lemma:Jordan块的初等因子}
    设矩阵\[
    \bm{A}=\mathrm{diag}\,\left\{
        \bm{J}_{r_1}\left(\lambda_1\right),\bm{J}_{r_2}\left(\lambda_2\right),\cdots,\bm{J}_{r_k}\left(\lambda_k\right)
    \right\}
    \]则$\bm{A}$的初等因子为\[
    \left(\lambda-\lambda_1\right)^{r_1},\left(\lambda-\lambda_2\right)^{r_2},\cdots,\left(\lambda-\lambda_k\right)^{r_k}
    \]
\end{lemma}
\begin{Proof}
因为\[
\lambda\bm{I}-\bm{A}=\begin{pmatrix}
    \lambda\bm{I}-\bm{J}_{r_1}\left(\lambda_1\right)&&&\\
    &\lambda\bm{I}-\bm{J}_{r_2}\left(\lambda_2\right)&&\\
    &&\ddots&\\
    &&&\lambda\bm{I}-\bm{J}_{r_k}\left(\lambda_k\right)
\end{pmatrix}
\]由\cref{lemma:Jordan块}知$\lambda\bm{I}-\bm{J}_{r_i}\left(\lambda_i\right)$相抵于\[
\mathrm{diag}\,\left\{
    1,\cdots,1,\left(\lambda-\lambda_i\right)^{r_i}
\right\}
\]即$\lambda\bm{I}-\bm{A}$相抵于\[
\mathrm{diag}\,\left\{
    1,\cdots,1,\left(\lambda-\lambda_1\right)^{r_1};1,\cdots,1,\left(\lambda-\lambda_2\right)^{r_2};\cdots;1,\cdots,1,\left(\lambda-\lambda_k\right)^{r_k}
\right\}
\]根据\cref{lemma:初等因子求得},$\bm{A}$的初等因子组为\[
\left(\lambda-\lambda_1\right)^{r_1},\left(\lambda-\lambda_2\right)^{r_2},\cdots,\left(\lambda-\lambda_k\right)^{r_k}\qedhere
\]
\end{Proof}
\end{green}
\begin{formal}
\begin{theorem}[Jordan标准型]\label{thm:Jordan标准型}
    设$\bm{A}\in M_n\left(\mathbb{C}\right)$的初等因子组为\[
    \left(\lambda-\lambda_1\right)^{r_1},\left(\lambda-\lambda_2\right)^{r_2},\cdots,\left(\lambda-\lambda_k\right)^{r_k}
    \]则$\bm{A}$相似于\[
    \bm{J}= \mathrm{diag}\,\left\{
        \bm{J}_{r_1}\left(\lambda_1\right),\bm{J}_{r_2}\left(\lambda_2\right),\cdots,\bm{J}_{r_k}\left(\lambda_k\right)
    \right\}
        \]
\end{theorem}
\begin{Proof}
由\cref{lemma:Jordan块的初等因子}知$\bm{J}$的初等因子组为\[
\left(\lambda-\lambda_1\right)^{r_1},\left(\lambda-\lambda_2\right)^{r_2},\cdots,\left(\lambda-\lambda_k\right)^{r_k}
\]由\cref{thm:全系不变量}知初等因子组是相似关系下的全系不变量,故$\bm{A}$相似于$\bm{J}$.
\end{Proof}
\end{formal}
\begin{red}
    \begin{remark}
        \cref{thm:Jordan标准型}需要保证在复数域$\mathbb{C}$上.
    \end{remark}
\end{red}
\begin{formal}
\begin{definition}[Jordan标准型]\label{def:Jordan标准型}
    \cref{thm:Jordan标准型}中的\[
    \bm{J}= \mathrm{diag}\,\left\{
        \bm{J}_{r_1}\left(\lambda_1\right),\bm{J}_{r_2}\left(\lambda_2\right),\cdots,\bm{J}_{r_k}\left(\lambda_k\right)
    \right\}
    \]称为$\bm{A}$的\textup{Jordan}标准型,其中$\bm{J}_{r_i}\left(\lambda_i\right)$称为$\bm{A}$的一个\textup{Jordan}块.
\end{definition}
\end{formal}
\begin{red}
\begin{remark}
    \textup{Jordan}标准型中\textup{Jordan}块的次序是无所谓的.但一个\textup{Jordan}块被一个初等因子完全唯一地决定.从而在不考虑\textup{Jordan}块的次序时,\textup{Jordan}标准型由矩阵$\bm{A}$完全唯一地决定.
\end{remark}
\end{red}
\begin{formal}
\begin{theorem}[几何意义的Jordan标准型]\label{thm:几何意义的Jordan标准型}
    设$\bm{\varphi}\in \mathcal{L}\left(
        V_{\mathbb{C}}^n
    \right)$,则存在$V$的一组基$\left\{
        \bm{e}_1,\bm{e}_2,\cdots,\bm{e}_n
    \right\}$使得$\bm{\varphi}$在这组基下的表示矩阵为\cref{thm:Jordan标准型}中的\textup{Jordan}标准型即\[
    \bm{J}= \mathrm{diag}\,\left\{
        \bm{J}_{r_1}\left(\lambda_1\right),\bm{J}_{r_2}\left(\lambda_2\right),\cdots,\bm{J}_{r_k}\left(\lambda_k\right)
    \right\}
    \]
\end{theorem}
\end{formal}
\begin{green}
\begin{corollary}[可对角化等价条件]\label{cor:可对角化等价条件}
    设$\bm{A}\in M_n\left(\mathbb{C}\right)$,则下列结论等价：

    \begin{enumerate}[label={\textup{(\arabic*)}}]
        \item $\bm{A}$可对角化
        \item $ \bm{A}$的极小多项式无重根
        \item $\bm{A}$的初等因子全为一次多项式（或者\textup{Jordan}块都是一阶的）
    \end{enumerate}
\end{corollary}
\begin{Proof}
    $(1)\Longrightarrow (2)$根据\cref{ex:极小多项式无重根等价于可对角化}即得.

    $(2)\Longrightarrow (3)$设$\bm{A}$的不变因子组为\[
    1,\cdots,1,d_1\left(\lambda\right),\cdots,d_k\left(\lambda\right)
    \]由\cref{thm:极小多项式}知$d_k\left(\lambda\right)$为$\bm{A}$的极小多项式$m\left(\lambda\right)$,从而$d_k\left(\lambda\right)$无重根,由整除关系知$d_1\left(\lambda\right),d_2\left(\lambda\right),\cdots,d_k\left(\lambda\right)$均没有重根,即它们的准素因子全体为一次多项式,从而初等因子全为一次多项式.

    $(3)\Longrightarrow (1)$设$\bm{A}$的初等因子组为\[
    \lambda-\lambda_1,\lambda-\lambda_2,\cdots,\lambda-\lambda_n
    \]于是Jordan标准型为\[
    \bm{J}=\begin{pmatrix}
        \lambda_1&&&\\
        &\lambda_2&&\\
        &&\ddots&\\
        &&&\lambda_n
    \end{pmatrix}
    \]即$\bm{A}$可对角化.
\end{Proof}
\end{green}
\begin{green}
\begin{corollary}[限制可对角化]\label{cor:限制可对角化}
    设$\bm{\varphi}\in \mathcal{L}\left(
        V_{\mathbb{C}}^n
    \right)$,$V_0$是$\bm{\varphi}$-不变子空间,若$\bm{\varphi}$可对角化,则$\displaystyle \bm{\varphi}\Big|_{V_0}$可对角化.
\end{corollary}
\begin{Proof}
    设$\bm{\varphi},\bm{\varphi}\Big|_{V_0}$的极小多项式分别为$m\left(\lambda\right),g\left(\lambda\right).$考虑到\begin{align*}
        m\left(\bm{\varphi}\Big|_{V_0}\right)=m\left(\bm{\varphi}\right)\Big|_{V_0}=\bm{0}\Big|_{V_0}=\bm{0}
    \end{align*}于是极小多项式$g\left(\lambda\right)\mid m\left(\lambda\right).$

    若$\bm{\varphi}$可对角化,则由\cref{cor:可对角化等价条件}知极小多项式$m\left(\lambda\right)$没有重根,那么$g\left(\lambda\right)$也无重根,根据\cref{cor:可对角化等价条件}知$\bm{\varphi}\Big|_{V_0}$可对角化.
\end{Proof}
\end{green}
\begin{green}
\begin{corollary}[代数意义的限制可对角化]\label{cor:代数意义的限制可对角化}
    若\[
    \bm{M}=\begin{pmatrix}
        \bm{A}^m&\bm{C}\\
        \bm{O}&\bm{B}^n
    \end{pmatrix}
    \]可对角化,则$\bm{A},\bm{B}$可对角化.
\end{corollary}
\begin{Proof}
    设$\bm{M}$的极小多项式为$m\left(\lambda\right)$,则\begin{align*}
        \bm{O}=m\left(\bm{M}\right)=\begin{pmatrix}
            m\left(\bm{A}\right)&*\\
            \bm{O}&m\left(\bm{B}\right) 
        \end{pmatrix}
    \end{align*}从而$m\left(\bm{A}\right)=m\left(\bm{B}\right)=\bm{O}$即$m_{\bm{A}}\left(\lambda\right)\mid m\left(\lambda\right),m_{\bm{B}}\left(\lambda\right)\mid m\left(\lambda\right)$,因为由\cref{cor:可对角化等价条件}知$m\left(\lambda\right)$无重根,则$m_{\bm{A}}\left(\lambda\right),m_{\bm{B}}\left(\lambda\right)$无重根,于是由\cref{cor:可对角化等价条件}知$\bm{A},\bm{B}$可对角化.
\end{Proof}
\end{green}
\begin{green}
\begin{corollary}[多个限制可对角化]\label{cor:多个限制可对角化}
    设$\bm{\varphi}\in \mathcal{L}\left(
        V_{\mathbb{C}}^n
    \right),V=V_1\oplus V_2\oplus\cdots\oplus V_k$,其中$V_i$是$\bm{\varphi}$-不变子空间$,\forall 1\leqslant i\leqslant k.$则$\bm{\varphi}$可对角化当且仅当$\bm{\varphi}\Big|_{V_i}$可对角化$\left(\forall 1\leqslant i\leqslant k\right).$
\end{corollary}
\begin{Proof}
必要性由\cref{cor:代数意义的限制可对角化}即得.

考虑充分性,因为存在$V_i$的一组基$\left\{
    \bm{e}_{i1},\bm{e}_{i2},\cdots,\bm{e}_{in_i}
\right\}$使得$\bm{\varphi}\Big|_{V_i}$在这组基下的表示矩阵为对角阵.根据\cref{thm:直和的判定}知子空间的基可以合并为$V$的一组基,从而$\bm{\varphi}$在这组基下的表示矩阵为对角阵,即$\bm{\varphi}$可对角化.
\end{Proof}
\end{green}
\begin{green}
\begin{corollary}[代数意义的多个限制可对角化]\label{cor:代数意义的多个限制可对角化}
    设
    \[
    \bm{A}=\begin{pmatrix}
        \bm{A}_1&&&\\
        &\bm{A}_2&&\\
        &&\ddots&\\
        &&&\bm{A}_k
    \end{pmatrix}
    \]则$\bm{A}$可对角化当且仅当$\bm{A}_1,\bm{A}_2,\cdots,\bm{A}_k$均可对角化.
\end{corollary}
\begin{Proof}
    由\cref{prop:分块对角矩阵的极小多项式}知\[
    m\left(\lambda\right)=\left[
        m_{\bm{A}_1}\left(\lambda\right),m_{\bm{A}_2}\left(\lambda\right),\cdots,m_{\bm{A}_k}\left(\lambda\right)
    \right]
    \]然后根据\cref{cor:可对角化等价条件}即得.
\end{Proof}
\end{green}
\begin{green}
\begin{corollary}[较小的域]\label{cor:较小的域}
    设$\bm{A}\in M_n\left(\mathbb{K}\right)$的特征值全在$\mathbb{K}$中,则$\bm{A}$在$\mathbb{K}$上相似于其\textup{Jordan}标准型.
\end{corollary}
\begin{Proof}
    首先,$\bm{A}$必定复相似于\[
    \bm{J}= \mathrm{diag}\,\left\{
        \bm{J}_{r_1}\left(\lambda_1\right),\bm{J}_{r_2}\left(\lambda_2\right),\cdots,\bm{J}_{r_k}\left(\lambda_k\right)
    \right\}
    \]容易见得Jordan块均在$M_n\left(\mathbb{K}\right)$,根据相似关系在基域扩张下的不变性知$\bm{A}$在$\mathbb{K}$上相似于其Jordan标准型$\bm{J}$.
\end{Proof}
\end{green}
\begin{brown}
\begin{example}
    设$\bm{A}$的初等因子组为\[
    \lambda-1,\left(\lambda-1\right)^3,\left(\lambda+2\right)^2,\left(\lambda+2\right)^3\]求其\textup{Jordan}标准型.
\end{example}
\begin{solution}
    $\bm{A}$的\textup{Jordan}标准型为\[
    \mathrm{diag}\,\left\{
        1,\begin{pmatrix}
            1&1&\\
            &1&1\\
            &&1
        \end{pmatrix},\begin{pmatrix}
            -2&1\\
            &-2
        \end{pmatrix},
        \begin{pmatrix}
            -2&1&\\
            &-2&1\\
            &&-2
        \end{pmatrix}
    \right\}
    \]
\end{solution}
\end{brown}
\begin{brown}
\begin{example}
    设$\bm{\varphi}\in \mathcal{L}\left(V_{\mathbb{C}}^4\right)$,其在$\left\{
        \bm{e}_1,\bm{e}_2,\bm{e}_3,\bm{e}_4
    \right\}$的表示阵为\[\bm{A}=
    \begin{pmatrix}
        3&1&0&0\\
        -4&-1&0&0\\
        6&1&2&1\\
        -14&-5&-1&0
    \end{pmatrix}
    \]试求另一组基$\left\{
        \bm{f}_1,\bm{f}_2,\bm{f}_3,\bm{f}_4
    \right\}$使得$\bm{\varphi}$在这组基下的表示矩阵为\textup{Jordan}标准型并求从$\bm{e}_i$到$\bm{f}_i$的过渡矩阵.
\end{example}
\begin{solution}
    根据\cref{thm:不同基下的表示矩阵}知过渡矩阵$\bm{P}$使得\[
    \bm{P}^{-1}\bm{A}\bm{P}=\bm{J}\Longrightarrow
    \bm{AP}= \bm{PJ}
    \]因为$\bm{A}$的初等因子组为\[
    \left(\lambda-1\right)^2,\left(\lambda-1\right)^2
    \]则\[
    \bm{J}=\begin{pmatrix}
        1&1&&&\\
        &1 &&&\\
        &&1&1&\\
        &&&1
    \end{pmatrix}
    \]取$\bm{P}=\left(\bm{\alpha}_1,\bm{\alpha}_2,\bm{\alpha}_3,\bm{\alpha}_4\right)$代入\[
    \bm{AP}= \bm{PJ}
    \]得到\[
    \begin{cases*}
    \bm{A\alpha}_1=\bm{\alpha}_1,\\
    \bm{A\alpha}_2=\bm{\alpha}_1+\bm{\alpha}_2,\\
    \bm{A\alpha}_3=\bm{\alpha}_3,\\
    \bm{A\alpha}_4=\bm{\alpha}_3+\bm{\alpha}_4
    \end{cases*}
    \]即$\bm{\alpha}_1,\bm{\alpha}_3$是$\bm{A}$的特征值$1$的两个线性无关的特征向量,又\[
    \left(\bm{A}-\bm{I}\right)\bm{\alpha}_2=\bm{\alpha}_1,\left(\bm{A}-\bm{I}\right)\bm{\alpha}_4=\bm{\alpha}_3
    \]于是需要先求出$\bm{A}$的特征值$1$的两个线性无关的特征向量$\bm{\alpha}_1,\bm{\alpha}_3$,然后代入上述方程求出$\bm{\alpha}_2,\bm{\alpha}_4$.最终得到\[\bm{\alpha}_2=\begin{pmatrix}
        \cfrac{5}{4}\\-\cfrac{7}{2}\\0\\0
    \end{pmatrix}\Longleftarrow
    \bm{\alpha}_1=\begin{pmatrix}
        -1\\2\\4\\0
    \end{pmatrix},\bm{\alpha}_3=\begin{pmatrix}
        -1\\2\\0\\4
    \end{pmatrix}\Longrightarrow\bm{\alpha}_4=\begin{pmatrix}
        \cfrac{1}{4}\\-\cfrac{3}{2}\\0\\0
    \end{pmatrix}
    \]
即过渡矩阵\[
\bm{P}=\begin{pmatrix}
    -1&\cfrac{5}{4}&-1&\cfrac{1}{4}\\
    2&-\cfrac{7}{2}&2&-\cfrac{3}{2}\\
    4&0&0&0\\
    0&0&4&0
\end{pmatrix}
\]而所求基为\[
\left\{
    \bm{f}_1,\bm{f}_2,\bm{f}_3,\bm{f}_4
\right\}=\left\{
    -\bm{e}_1+2\bm{e}_2+4\bm{e}_3,\cfrac{5}{4}\bm{e}_1-\cfrac{7}{2}\bm{e}_2,-\bm{e}_1+2\bm{e}_2+4\bm{e}_4,\cfrac{1}{4}\bm{e}_1-\cfrac{3}{2}\bm{e}_2
\right\}
\]
\end{solution}
\end{brown}
\subsection{Jordan标准型的相关应用}
这一部分,首先将讨论$\bm{\varphi}$的特征值的度数（几何重数）和重数（代数重数）在Jordan标准型中的体现,然后给出全空间关于Jordan标准型的两种直和分解,最后给出Jordan标准型的一些应用.
\begin{red}
\begin{remark}
    总设$n$维线性空间$V$,线性变换$\bm{\varphi}\in \mathcal{L}\left(V\right),$其初等因子组为\[
    \left(\lambda-\lambda_1\right)^{r_1},\left(\lambda-\lambda_2\right)^{r_2},\cdots,\left(\lambda-\lambda_k\right)^{r_k}
    \]并且存在一组基\[
    \left\{
        e_{11},e_{12},\cdots,e_{1r_1};e_{21},e_{22},\cdots,e_{2r_2};\cdots;e_{k1},e_{k2},\cdots,e_{kr_k}
    \right\}
    \]使得$\bm{\varphi}$在这组基下的表示矩阵为\[
    \bm{J}= \mathrm{diag}\,\left\{
        \bm{J}_{r_1}\left(\lambda_1\right),\bm{J}_{r_2}\left(\lambda_2\right),\cdots,\bm{J}_{r_k}\left(\lambda_k\right)
    \right\}
    \]
\end{remark}
\end{red}因为$\forall 1\leqslant i\leqslant k$\[
    \left(
        \bm{\varphi}\left(\bm{e}_{i1}\right),\bm{\varphi}\left(\bm{e}_{i2}\right),\cdots,\bm{\varphi}\left(\bm{e}_{ir_i}\right)
    \right)=\left(
        \bm{e}_{i1},\bm{e}_{i2},\cdots,\bm{e}_{ir_i}
    \right)\begin{pmatrix}
        \lambda_i&1&&&\\
        &\lambda_i&1&\\
        &&\ddots&\ddots&\\
        &&&\ddots&1\\
        &&&&\lambda_i
    \end{pmatrix}
    \]即\begin{align*}
        \begin{cases*}
        \bm{\varphi}\left(\bm{e}_{i1}\right)=\lambda_i\bm{e}_{i1}\\
        \bm{\varphi}\left(\bm{e}_{i2}\right)=\bm{e}_{i1}+\lambda_i\bm{e}_{i2}\\
        \qquad\qquad\cdots\cdots\\
        \bm{\varphi}\left(\bm{e}_{ir_i}\right)=\bm{e}_{i,r_i-1}+\lambda_i\bm{e}_{ir_i}
        \end{cases*}\tag{*}
    \end{align*}令张成的子空间\[
    V_i=L\left(
        \bm{e}_{i1},\bm{e}_{i2},\cdots,\bm{e}_{ir_i}
    \right)
    \]根据上文,$\bm{\varphi}\left(V_i\right)\subseteq V_i$即$V_i$是$\bm{\varphi}$-不变子空间.根据\cref{thm:直和的判定},$V=V_1\oplus V_2\oplus\cdots\oplus V_k.$
    \begin{formal}
    \begin{definition}[Jordan块相伴的子空间]\label{def:Jordan块相伴的子空间}
        如上文,张成的子空间\[
        V_i=L\left(
            \bm{e}_{i1},\bm{e}_{i2},\cdots,\bm{e}_{ir_i}
        \right)
        \]称为\textup{Jordan}块$\bm{J}_{r_i}\left(\lambda_i\right)$相伴的子空间.
    \end{definition}
    \end{formal}

    设\[
    f\left(\lambda\right)=\left(\lambda-\lambda_1\right)^{r_1}\left(\lambda-\lambda_2\right)^{r_2}\cdots\left(\lambda-\lambda_k\right)^{r_k}
    \]不妨设$\lambda_1=\lambda_2=\cdots=\lambda_s\neq \lambda_j,\forall s<j\leqslant k.$于是\[
    f\left(\lambda\right)=\left(\lambda-\lambda_1\right)^{r_1+r_2+\cdots+r_s}\left(\lambda-\lambda_{s+1}\right)^{r_{s+1}}\cdots\left(\lambda-\lambda_k\right)^{r_k}
    \]其中$\lambda_1$的代数重数为$r_1+r_2+\cdots+r_s$,其等于属于特征值$\lambda_1$的Jordan块的阶数之和,而度数为\[
    \dim\,V_{\lambda_1}=\dim\,\mathrm{Ker}\,\left(
        \bm{\varphi}-\lambda_1\bm{I}_V
    \right)=n-\dim\,\mathrm{Im}\,\left(
        \bm{\varphi}-\lambda_1\bm{I}_V
    \right)=n-\mathrm{rank}\,\left(
        \bm{\varphi}-\lambda_1\bm{I}_V
    \right)
    \]考虑\[\bm{J}_{r_i}\left(\lambda_i\right)-\lambda_1\bm{I}=
    \begin{pmatrix}
        \lambda_i-\lambda_1&1&&&\\
        &\lambda_i-\lambda_1&1&\\
        &&\ddots&\ddots&\\
        &&&\ddots&1\\
        &&&&\lambda_i-\lambda_1
    \end{pmatrix}
    \]则秩\[
    \mathrm{rank}\,\left(
        \bm{J}_{r_i}\left(\lambda_i\right)-\lambda_1\bm{I}
    \right)=\begin{cases*}
    r_i-1&,$\lambda_i=\lambda_1\left(1\leqslant i\leqslant s\right)$\\
    r_i&,$\lambda_i\neq \lambda_1\left(s<i\leqslant k\right)$
    \end{cases*}
    \]于是\begin{align*}
    \mathrm{r}\,\left(\bm{J}-\lambda_1\bm{I}\right)&=\sum_{i=1}^k\mathrm{r}\,\left(
        \bm{J}_{r_i}\left(\lambda_i\right)-\lambda_1\bm{I}
    \right)\\
    &=r_1-1+r_2-1+\cdots+r_s-1+r_{s+1}+\cdots+r_k\\
    &=n-s
    \end{align*}
    于是$\lambda_1$的几何重数为$s$即属于特征值$\lambda_1$的Jordan块的个数.
\begin{formal}
\begin{theorem}[代数重数与几何重数]\label{thm:代数重数与几何重数}
    特征值$\lambda_1$的代数重数等于属于特征值$\lambda_1$的所有\textup{Jordan}块的阶数之和,而几何重数等于属于特征值$\lambda_1$的\textup{Jordan}块的个数.
\end{theorem}
\end{formal}
\begin{red}
\begin{remark}
    $*$式中$\bm{e}_{i1}$是特征向量,而$\bm{e}_{i2},\bm{e}_{i3},\cdots,\bm{e}_{ir_i}$被称为广义特征向量.
\end{remark}
\end{red}
\begin{green}
\begin{corollary}[特征子空间的基]\label{cor:特征子空间的基}
    $\left\{
        \bm{e}_{11},\bm{e}_{21},\cdots,\bm{e}_{s1}
    \right\}$是特征子空间$V_{\lambda_1}$的一组基.
\end{corollary}
\end{green}
令$\bm{\psi}=\bm{\varphi}-\lambda_i\bm{I}_V$,则$*$式成为\[
\begin{cases*}
    \bm{\psi}\left(\bm{e}_{i1}\right)=\bm{0}\\
    \bm{\psi}\left(\bm{e}_{i2}\right)=\bm{e}_{i1}\\
    \qquad\cdots\cdots\\
    \bm{\psi}\left(\bm{e}_{ir_i}\right)=\bm{e}_{i,r_i-1}
\end{cases*}
\]即有\[
\bm{e}_{ir_i}\xlongrightarrow{\bm{\psi}}\bm{e}_{i,r_i-1}\xlongrightarrow{\bm{\psi}}\cdots\xlongrightarrow{\bm{\psi}}\bm{e}_{i1}\xlongrightarrow{\bm{\psi}}\bm{0}
\]
\begin{formal}
\begin{definition}[循环子空间]\label{def:循环子空间}
    设$\bm{\psi}\in \mathcal{L}\left(V\right)$,$V_0$是$r$维的$\bm{\psi}$-不变子空间,若存在$\bm{0}\neq \bm{\alpha}\in V_0\,\textup{s.t.}$\[
    \left\{
        \bm{\alpha},\bm{\psi}\left(\bm{\alpha}\right),\cdots,\bm{\psi}^{r-1}\left(\bm{\alpha}\right)
    \right\}
    \]是$V_0$的一组基,则称$V_0$是$\bm{\psi}$的一个循环子空间.$\bm{\alpha}$称为该循环子空间的循环向量.
\end{definition}
\end{formal}
\begin{red}
    \begin{remark}
        定义一般的循环子空间时不需要保证$\bm{\psi}^r\left(\bm{\alpha}\right)=\bm{0}.$
    \end{remark}
\end{red}
于是对于上式\[
\left\{
    \bm{e}_{ir_i},\bm{\psi}\left(\bm{e}_{ir_i}\right),\cdots,\bm{\psi}^{r_i-1}\left(\bm{e}_{ir_i}\right)
\right\}
\]是$V_i$的一组基,于是$V_i$是$\bm{\psi}=\bm{\varphi}-\lambda_i\bm{I}_V$的一个循环子空间,而$\bm{e}_{ir_i}$是$V_i$的循环向量$\left(1\leqslant i\leqslant k\right)$.

于是全空间一定可以分解为循环子空间的直和.
\begin{green}
\begin{lemma}[根子空间]\label{lemma:根子空间}
    设\[
    R\left(\lambda_1\right)=V_1\oplus V_2\oplus\cdots\oplus V_s
    \]则\begin{align*}
    R\left(\lambda_1\right)&=\mathrm{Ker}\,\left(
        \bm{\varphi}-\lambda_1\bm{I}_V
    \right)^n\\
    &=\left\{
        \bm{v}\in V\mid \left(
            \bm{\varphi}-\lambda_1\bm{I}_V
            \right)^n\left(\bm{v}\right)=\bm{0}
    \right\}
    \end{align*}
\end{lemma}
\begin{Proof}
    先证$R\left(\lambda_1\right)\subseteq\mathrm{Ker}\,\left(\bm{\varphi}-\lambda_1\bm{I}_V\right)^n$.任取\[
    R\left(\lambda_1\right)\ni\bm{v}=\bm{v}_1+\bm{v}_2+\cdots+\bm{v}_s,\bm{v}_i\in V_i
    \]考虑$V_i$的基\[
    \bm{e}_{ir_i}\xlongrightarrow{\bm{\varphi}-\lambda_1\bm{I}_V}\bm{e}_{i,r_i-1}\xlongrightarrow{\bm{\varphi}-\lambda_1\bm{I}_V}\cdots\xlongrightarrow{\bm{\varphi}-\lambda_1\bm{I}_V}\bm{e}_{i1}\xlongrightarrow{\bm{\varphi}-\lambda_1\bm{I}_V}\bm{0}
    \]即$\left(
        \bm{\varphi}-\lambda_1\bm{I}_V
    \right)^{r_i}\left(
        \bm{e}_{ij}
    \right)=\bm{0},\forall j\Longrightarrow
    \left(
        \bm{\varphi}-\lambda_1\bm{I}_V
    \right)^{r_i}\left(
        \bm{v}_i
    \right)=\bm{0}$.于是\[
    \left(
        \bm{\varphi}-\lambda_1\bm{I}_V
    \right)^n\left(\bm{v}\right)=
    \left(
        \bm{\varphi}-\lambda_1\bm{I}_V
    \right)^n\left(\bm{v}_1\right)+\left(
        \bm{\varphi}-\lambda_1\bm{I}_V
    \right)^n\left(\bm{v}_2\right)+\cdots+\left(
        \bm{\varphi}-\lambda_1\bm{I}_V
    \right)^n\left(\bm{v}_s\right)=\bm{0}
    \]

    再证$\mathrm{Ker}\,\left(\bm{\varphi}-\lambda_1\bm{I}_V\right)^n\subseteq R\left(\lambda_1\right)$.设$\bm{v}\in V$在基下的坐标向量\[
    \bm{x}=\left(
        x_{11},x_{12},\cdots,x_{1r_1},x_{21},x_{22},\cdots,x_{2r_2},\cdots,x_{s1},x_{s2},\cdots,x_{sr_s}
    \right)'
    \]考虑$\left(
        \bm{J}-\lambda_1\bm{I}
    \right)^n\bm{x}=\bm{0}$的解.

    由上文知\[
    \left(
        \bm{J}_{r_i}\left(\lambda_i\right)-\lambda_1\bm{I}
    \right)^n=\begin{cases*}
        \bm{0}&,$1\leqslant i\leqslant s$\\
        \text{非异}&,$s<i\leqslant k$
    \end{cases*}
    \]于是\[
    x_{i1}=x_{i2}=\cdots=x_{ir_i}=0,\forall s<i\leqslant k
    \]而\[
    \left(
        x_{11},x_{12},\cdots,x_{1r_1},x_{21},x_{22},\cdots,x_{2r_2},\cdots,x_{s1},x_{s2},\cdots,x_{sr_s}
    \right)
    \]为任意解.于是\[
    \mathrm{Ker}\,\left(
        \bm{\varphi}-\lambda_1\bm{I}_V
    \right)^n=V_1\oplus V_2\oplus\cdots\oplus V_s=R\left(\lambda_1\right)\qedhere
    \]
\end{Proof}
\end{green}
\begin{red}
    \begin{remark}
        事实上,\cref{lemma:根子空间}的证明中用到的系数$n$最佳取到$\max\,\left\{
            r_1,r_2,\cdots,r_s
        \right\}$或者退一步取到$r_1+r_2+\cdots+r_s.$
    \end{remark}
\end{red}
\begin{formal}
\begin{definition}[根子空间]\label{def:根子空间}
    \[
    R\left(\lambda_1\right)=\mathrm{Ker}\,\left(
        \bm{\varphi}-\lambda_1\bm{I}_V
    \right)^n=V_1\oplus V_2\oplus\cdots\oplus V_s
    \]称为特征值$\lambda_1$的根子空间.
\end{definition}
\end{formal}
\begin{red}
\begin{remark}
    同上,\cref{def:根子空间}中的$n$最佳取到$\max\,\left\{
        r_1,r_2,\cdots,r_s
    \right\}$或者退一步取到$r_1+r_2+\cdots+r_s$即代数重数（具体在\cref{ex:像空间与核空间稳定}中已证明）.
\end{remark}
\end{red}
\begin{red}
\begin{remark}
    根据\cref{lemma:幂的像空间与核空间},特征值$\lambda_1$的特征子空间$V_{\lambda_1}=\mathrm{Ker}\,\left(\bm{\varphi}-\lambda_1\bm{I}_V\right)\subseteq R\left(\lambda_1\right).$
\end{remark}
\end{red}
\begin{green}
\begin{corollary}[可对角化条件]\label{cor:可对角化条件}
    $\bm{\varphi}\in \mathcal{L}\left(V_{\mathbb{C}}\right)$可对角化当且仅当对于所有特征值$\lambda_0$均有\[
    V_{\lambda_0}=R\left(\lambda_0\right)
    \]
\end{corollary}
\begin{Proof}
    考虑维数,根据\cref{thm:特征子空间的维数},左侧特征子空间的维数是$\bm{\varphi}$的几何重数；右侧根据定义是各循环子空间的维数之和,而循环子空间的维数之和为属于$\lambda_1$的\textup{Jordan}块的阶数之和,即代数重数.于是根据\cref{def:完全特征向量系}和\cref{thm:可对角化的条件2}知$\bm{\varphi}$可对角化.
\end{Proof}
\end{green}
下面对上面叙述的内容进行一个总结.
\begin{formal}
\begin{theorem}[全空间关于Jordan标准型的两种直和分解]\label{thm:全空间关于Jordan标准型的两种直和分解}
    设$\bm{\varphi}\in \mathcal{L}\left(V_{\mathbb{C}}^n\right)$\begin{enumerate}[label={\textup{(\arabic*)}}]
        \item 设$\bm{\varphi}$的初等因子组为\[
        \left(
            \lambda-\lambda_1
        \right)^{r_1},\left(
            \lambda-\lambda_2
        \right)^{r_2},\cdots,\left(
            \lambda-\lambda_k
        \right)^{r_k}  
        \]并且它们相伴的子空间为$V_1,V_2,\cdots,V_k$,则\[
        V=V_1\oplus V_2\oplus\cdots\oplus V_k
        \]其中$V_i$是维数等于$r_i$的关于$\bm{\varphi}-\lambda_i\bm{I}_V$的循环子空间.
        \item 设$\bm{\varphi}$的全体不同特征值为$\lambda_1,\lambda_2,\cdots,\lambda_s$,则\[
        V=R\left(\lambda_1\right)\oplus R\left(\lambda_2\right)\oplus\cdots\oplus R\left(\lambda_s\right)
        \]其中$R\left(\lambda_i\right)$是特征值$\lambda_i$的根子空间,其维数等于$\lambda_i$的代数重数,并且每一个$R\left(\lambda_i\right)$都可以分解为$(1)$中若干个$V_j$的直和.
    \end{enumerate}
\end{theorem}
\end{formal}
\begin{red}
    \begin{remark}
        \textup{Jordan}标准型理论的核心是如果一个矩阵问题的条件与结论在相似关系下不发生改变,那么就可以将该问题化约到\textup{Jordan}标准型的情况下进行讨论,并且进一步地化约到\textup{Jordan}块的情况下进行讨论.

        但在具体实践中,一般来说是\[
        \text{\textup{Jordan}块成立}\Longrightarrow
        \text{\textup{Jordan}标准型成立}\Longrightarrow
        \text{一般矩阵成立}
        \]
    \end{remark}
\end{red}
\begin{brown}
    \begin{example}
        设$\bm{A}\in M_n\left(\mathbb{C}\right)$,证明：$\bm{A}=\bm{BC},$其中$\bm{B},\bm{C}$为对称阵.
    \end{example}
    \begin{Proof}
        因为存在非异阵$\bm{P}$使得\[
        \bm{P}^{-1}\bm{A}\bm{P}=\bm{J}=
        \mathrm{diag}\,\left\{
            \bm{J}_{r_1}\left(\lambda_1\right),\bm{J}_{r_2}\left(\lambda_2\right),\cdots,\bm{J}_{r_k}\left(\lambda_k\right)
        \right\}
        \]先来考虑Jordan块\[
        \bm{J}_{r_i}\left(\lambda_i\right)=\begin{pmatrix}
            \lambda_i&1&&&\\
            &\lambda_i&1&\\
            &&\ddots&\ddots&\\
            &&&\ddots&1\\
            &&&&\lambda_i
        \end{pmatrix}
        \]事实上,这个矩阵旋转一定角度后自然成为一个对称阵,注意到\[
        \begin{pmatrix}
            \lambda_i&1&&&\\
            &\lambda_i&1&\\
            &&\ddots&\ddots&\\
            &&&\ddots&1\\
            &&&&\lambda_i
        \end{pmatrix}=\begin{pmatrix}
            &&1&\lambda_i\\
            &1&\lambda_i\\
            &\iddots&\iddots&\\
            1&\iddots\\
            \lambda_i
        \end{pmatrix}\begin{pmatrix}
            &&&&1\\
            &&&1\\
            &&\iddots\\
            &1\\
            1
        \end{pmatrix}=\bm{S}_i\bm{T}_i
        \]然后推广到Jordan标准型即令\begin{align*}
           & \bm{S}=\mathrm{diag}\,\left\{
                \bm{S}_1,\bm{S}_2,\cdots,\bm{S}_k
            \right\}\\
            & \bm{T}=\mathrm{diag}\,\left\{
                \bm{T}_1,\bm{T}_2,\cdots,\bm{T}_k
            \right\}
        \end{align*}于是\[
        \bm{J}=\bm{ST}
        \]其中$\bm{S},\bm{T}$是对称阵.又因为$\bm{A}=\bm{PJP}^{-1}=\bm{PSTP}^{-1}=\left(\bm{PSP}'\right)\left(\left(\bm{P}^{-1}\right)'\bm{TP}^{-1}\right)$是对称阵.
    \end{Proof}
\end{brown}
\begin{brown}
    \begin{example}
        设$\bm{A}\in M_n\left(\mathbb{C}\right)$,求$\bm{A}^k\left(k\geqslant 1\right).$
    \end{example}
    \begin{solution}
        因为存在非异阵$\bm{P}$使得\[
        \bm{P}^{-1}\bm{A}\bm{P}=\bm{J}=
        \mathrm{diag}\,\left\{
            \bm{J}_{r_1}\left(\lambda_1\right),\bm{J}_{r_2}\left(\lambda_2\right),\cdots,\bm{J}_{r_k}\left(\lambda_k\right)
        \right\}
        \]于是$\bm{A}^k=\bm{P}\bm{J}^k\bm{P}^{-1}=\bm{P}\mathrm{diag}\,\left\{
            \bm{J}_{r_1}\left(\lambda_1\right)^k,\bm{J}_{r_2}\left(\lambda_2\right)^k,\cdots,\bm{J}_{r_k}\left(\lambda_k\right)^k
        \right\}\bm{P}^{-1}$.

        考虑\begin{align*}\bm{J}_n\left(\lambda_0\right)&=
        \begin{pmatrix}
            \lambda_0&1&&&\\
            &\lambda_0&1&\\
            &&\ddots&\ddots&\\
            &&&\ddots&1\\
            &&&&\lambda_0
        \end{pmatrix}\\
        &=\lambda_0\bm{I}_n+\bm{N}
    \end{align*}其中\[
    \bm{N}=\bm{J}_n\left(0\right)=\begin{pmatrix}
        0&1&&&\\
        &0&1&\\
        &&\ddots&\ddots&\\
        &&&\ddots&1\\
        &&&&0
    \end{pmatrix}
    \]并且有良好的性质\[
    \bm{N}^2=\begin{pmatrix}
        &&1\\
        &&&1\\
        &&&&\ddots\\
        &&&&&1\\
        &&&&&&\\
        &&&&&&&
    \end{pmatrix},\cdots,\bm{N}^{n-1}=\begin{pmatrix}
        &&&&&&&1\\
        \\
        \\
        \\
        \\
        \\
    \end{pmatrix}
    \]因为$\bm{N}$的特征多项式为$\lambda^n$,于是根据\cref{thm:Cayley-Hamilton定理}\textup{Cayley-Hamilton}定理知$\bm{N}^n=\bm{O}.$因为纯量阵与任何矩阵都可交换,则设$k\geqslant n-1$有\begin{align*}
        \bm{J}_n\left(\lambda_0\right)^k&=\left(\lambda_0\bm{I}_n+\bm{N}\right)^k\\
        &=\lambda_0^k\bm{I}_n+C_k^1\lambda_0^{k-1}\bm{N}+C_k^2\lambda_0^{k-2}\bm{N}^2+\cdots+C_k^{n-1}\lambda_0^{k-n+1}\bm{N}^{n-1}\\
        &=\begin{pmatrix}
            \lambda_0^k&C_k^1\lambda_0^{k-1}&C_k^2\lambda_0^{k-2}&\cdots&C_k^{n-1}\lambda_0^{k-n+1}\\
            &\lambda_0^k&C_k^1\lambda_0^{k-1}&\cdots&C_k^{n-2}\lambda_0^{k-n+2}\\
            &&\lambda_0^k&\cdots&\vdots\\
            &&&\ddots&\vdots\\
            &&&&\lambda_0^k
        \end{pmatrix}
    \end{align*}合并即可.
    \end{solution}
\end{brown}
\begin{green}
\begin{lemma}[同时对角化]\label{lemma:同时对角化}
    设$\bm{A},\bm{B}\in M_n\left(\mathbb{C}\right)$,若$\bm{A},\bm{B}$均可对角化且乘积可交换即$\bm{AB}=\bm{BA}$,则$\bm{A},\bm{B}$可同时对角化即存在非异阵$\bm{P}$使得$\bm{P}^{-1}\bm{AP}$和$\bm{P}^{-1}\bm{BP}$均为对角阵.
\end{lemma}
\begin{Proof}
    考虑几何角度,即$\bm{\varphi},\bm{\psi}\in\mathcal{L}\left(V_{\mathbb{C}}^n\right)$,二者均可对角化且$\bm{\varphi}\bm{\psi}=\bm{\psi}\bm{\varphi}$,则$\bm{\varphi},\bm{\psi}$可同时对角化即存在一组基使得$\bm{\varphi},\bm{\psi}$在这组基下的表示矩阵均为对角阵.

    下面对线性空间的维数$n$进行归纳.$n=1$显然.设维数$<n$时成立,考虑维数等于$n$的情况.

    设$\bm{\varphi}$的全体不同特征值为$\lambda_1,\lambda_2,\cdots,\lambda_s$.一方面,$s=1$则$\bm{\varphi}=\lambda_1\bm{I}_V$.因为$\bm{\psi}$可对角化即存在一组基$\left\{
        \bm{e}_1,\bm{e}_2,\cdots,\bm{e}_n
    \right\}$使得$\bm{\psi}$在这组基下的表示矩阵为对角阵,因为纯量变换$\bm{\varphi}=\lambda_1\bm{I}_V$在任意基下的表示矩阵均为对角阵$\lambda_1\bm{I}_n$,于是证毕.

    另一方面,$s>1$时.设对应特征子空间$V_1,V_2,\cdots,V_s$.根据\cref{thm:可对角化的条件},$\bm{\varphi}$可对角化即\[
    V=V_1\oplus V_2\oplus\cdots\oplus V_s
    \]其中$\dim\,V_i<n,\forall 1\leqslant i\leqslant s.$

    下面利用$\bm{\varphi\psi}=\bm{\psi\varphi}$证明$\bm{\varphi}$的特征子空间$V_i$是$\bm{\psi}$的不变子空间.任取$\bm{v}_i\in V_i\Longrightarrow \bm{\varphi}\left(\bm{v}_i\right)=\lambda_i\bm{v}_i.$于是\begin{align*}
        \bm{\varphi}\left(\bm{\psi}\left(\bm{v}_i\right)\right)=\bm{\psi}\left(\bm{\varphi}\left(\bm{v}_i\right)\right)=\lambda_i\bm{\psi}\left(\bm{v}_i\right)
    \end{align*}
    于是$\bm{\psi}\left(\bm{v}_i\right)\in V_i$即$V_i$是$\bm{\psi}$的不变子空间.

    做限制$\bm{\varphi}\Big|_{V_i},\bm{\psi}\Big|_{V_i}$根据\cref{cor:多个限制可对角化}知它们均可对角化并且可交换.并且维数在$s>1$前提下严格小于$n$,则由归纳假设知存在$V_i$的一组基$\left\{
        \bm{e}_{i1},\bm{e}_{i2},\cdots,\bm{e}_{ir_i}
    \right\}$使得$\bm{\varphi}\Big|_{V_i},\bm{\psi}\Big|_{V_i}$在这组基下的表示矩阵均为对角阵.

    考虑到直和$V=V_1\oplus V_2\oplus\cdots\oplus V_s\Longrightarrow$子空间的基可以拼成全空间的一组基,那么$\bm{\varphi},\bm{\psi}$在这组基下的表示矩阵即由各限制拼成的对角阵.于是证毕.
\end{Proof}
\end{green}
\begin{formal}
\begin{theorem}[Jordan-Chevalley分解定理]\label{thm:Jordan-Chevalley分解定理}
    设$\bm{A}\in M_n\left(\mathbb{C}\right),$则$\bm{A}$一定可以作分解$\bm{A}=\bm{B}+\bm{C}$,其中\begin{enumerate}[label={\textup{(\arabic*)}}]
        \item $\bm{B}$可对角化
        \item $\bm{C}$是幂零阵
        \item 乘积可交换：$\bm{BC}=\bm{CB}$
        \item $\bm{B},\bm{C}$均可表示为$\bm{A}$的多项式
    \end{enumerate}并且满足$(1)\sim (3)$的分解是唯一的.
\end{theorem}
\begin{Proof}
    设非异阵$\bm{P}$使得\[
    \bm{P}^{-1}\bm{A}\bm{P}=\mathrm{diag}\,\left\{
        \bm{J}_1,\bm{J}_2,\cdots,\bm{J}_s
    \right\}
    \]其中$\lambda_1,\lambda_2,\cdots,\lambda_s$是$\bm{A}$的全体不同特征值,而$\bm{J}_i$是$\lambda_i$的根子空间的块（阶数为$\lambda_i$的代数重数$m_i$）\[
    \bm{J}_i=\begin{pmatrix}
        \bm{J}_{r_1}\left(\lambda_i\right)&&&\\
        &\bm{J}_{r_2}\left(\lambda_i\right)&&\\
        &&\ddots&\\
        &&&\bm{J}_{r_k}\left(\lambda_i\right)
    \end{pmatrix}=\lambda_i\bm{I}+\bm{N}_i=\bm{M}_i+\bm{N}_i
    \]前者是对角阵,后者每个分块的上次对角线元素为$1$并且是幂零阵（$\bm{N}_i^{m_i}=\bm{O}$）,并且有$\bm{M}_i\bm{N}_i=\bm{N}_i\bm{M}_i$.

    做拼接\begin{align*}
        &\bm{M}=\mathrm{diag}\,\left\{
            \bm{M}_1,\bm{M}_2,\cdots,\bm{M}_s
        \right\}\\
        &\bm{N}=\mathrm{diag}\,\left\{
            \bm{N}_1,\bm{N}_2,\cdots,\bm{N}_s
        \right\}
    \end{align*}前者是对角阵,后者是幂零阵并且可交换即$\bm{MN}=\bm{NM}$.于是分解$\bm{J}=\bm{M}+\bm{N}$满足条件$(1)\sim (3)$.

    $\bm{J}_i$的特征多项式为$\left(\lambda-\lambda_i\right)^{m_i}$,由\cref{thm:Cayley-Hamilton定理}知$
    \left(
        \bm{J}_i-\lambda_i\bm{I}
    \right)^{m_i}=
        \bm{N}_i^{m_i}=\bm{O}
    $.因为\[
    \left(\lambda-\lambda_1\right)^{m_1}\left(\lambda-\lambda_2\right)^{m_2}\cdots\left(\lambda-\lambda_s\right)^{m_s}
    \]两两互素,则由\cref{thm:Chinese Remainden Theorem}知存在$g\left(\lambda\right)$使得\[
    g\left(\lambda\right)=\left(\lambda-\lambda_i\right)^{m_i}q_i\left(\lambda\right)+\lambda_i,\forall 1\leqslant i\leqslant s
    \]于是\begin{align*}
        g\left(\bm{J}_i\right)=\lambda_i\bm{I}=\bm{M}_i
    \end{align*}而\begin{align*}
        g\left(\bm{J}\right)&=\mathrm{diag}\,\left\{
            g\left(\bm{J}_1\right),g\left(\bm{J}_2\right),\cdots,g\left(\bm{J}_s\right)
        \right\}\\
        &=\mathrm{diag}\,\left\{
            \bm{M}_1,\bm{M}_2,\cdots,\bm{M}_s
        \right\}=\bm{M}
    \end{align*}并且$\bm{N}=\bm{J}-g\left(\bm{J}\right)$即$(4)$证毕.

    再考虑到$\bm{A}=\bm{P}\bm{JP}^{-1}=\bm{P}\left(\bm{M}+\bm{N}\right)\bm{P}^{-1}=\bm{P}\bm{MP}^{-1}+\bm{P}\bm{NP}^{-1}=\bm{B}+\bm{C}$,容易验证$\bm{B}$可对角化,而$\bm{C}$是幂零阵并且\begin{align*}
        \bm{BC}=\bm{PMNP}^{-1}=\bm{PNMP}^{-1}=\bm{CB}
    \end{align*}同时\begin{align*}
        &g\left(\bm{A}\right)=\bm{P}g\left(\bm{J}\right)\bm{P}^{-1}=\bm{B}\\
        &\bm{A}-\bm{B}=\bm{A}-g\left(\bm{A}\right)=\bm{C}
    \end{align*}

    于是Jordan-Chevalley分解存在性证毕,下面证明唯一性.再设$\bm{A}=\bm{B}_1+\bm{C}_1,$其中$\bm{B}_1,\bm{C}_1$满足$(1)\sim (3)$,容易证明$\bm{B}_1,\bm{C}_1$与$\bm{A}$是可交换的,于是$\bm{B}_1$与$\bm{A}$的任何一个多项式可交换即与$\bm{B}$可交换即$\bm{BB}_1=\bm{B}_1\bm{B}$.同理$\bm{CC}_1=\bm{C}_1\bm{C}$.因为\[
    \bm{B}-\bm{B}_1=\bm{C}_1-\bm{C}
    \]不妨设$\bm{C}_1^r=\bm{O},\bm{C}^s=\bm{O}$则$\displaystyle \left(\bm{C}_1-\bm{C}\right)^{r+s}=\sum_{i+j=r+s}C_{r+s}^i\bm{C}_1^i\bm{C}^j=\bm{O}$于是$\left(
        \bm{B}-\bm{B}_1
    \right)^{r+s}=\bm{O}$.由\cref{lemma:同时对角化}知存在非异阵$\bm{P}$使得\[
    \bm{P}^{-1}\bm{BP}=\begin{pmatrix}
        a_1&&&\\
        &a_2&&\\
        &&\ddots&\\
        &&&a_n
    \end{pmatrix},\bm{P}^{-1}\bm{B}_1\bm{P}=\begin{pmatrix}
        b_1&&&\\
        &b_2&&\\
        &&\ddots&\\
        &&&b_n
    \end{pmatrix}
    \]于是\begin{align*}
    \bm{O}&=\bm{P}^{-1}\left(\bm{B}-\bm{B}_1\right)^{r+s}\bm{P}\\
    &=\left(
        \bm{P}^{-1}\bm{B}\bm{P}-\bm{P}^{-1}\bm{B}_1\bm{P}
    \right)^{r+s}\\
    &=\begin{pmatrix}
        \left(a_1-b_1\right)^{r+s}&&&\\
        &\left(a_2-b_2\right)^{r+s}&&\\
        &&\ddots&\\
        &&&\left(a_n-b_n\right)^{r+s}
    \end{pmatrix}
\end{align*}于是$a_1=b_1,a_2=b_2,\cdots,a_n=b_n$即$\bm{B}=\bm{B}_1\Longrightarrow \bm{C}=\bm{C}_1$.

    于是Jordan-Chevalley分解定理证毕.
\end{Proof}
\end{formal}
\subsection{Jordan标准型应用三大主题}
\begin{red}
    \begin{remark}
        \textup{Jordan}标准型应用的第一个主题是如何合理地选取特征向量求得广义特征向量
    \end{remark}
\end{red}
\begin{brown}
    \begin{example}
        考虑矩阵\[
        \bm{A}=\begin{pmatrix}
            2&6&-15\\1&1&-5\\1&2&-6
        \end{pmatrix}\]求可逆阵$\bm{P}$使得$\bm{P}^{-1}\bm{A}\bm{P}$为\textup{Jordan}标准型.
    \end{example}
    \begin{solution}
        计算得到初等因子组\[\lambda+1,
        \left(\lambda+1\right)^2
        \]即\textup{Jordan}标准型为\[
        \begin{pmatrix}
            -1&&\\
            &-1&1\\
            &&-1
        \end{pmatrix}
        \]设$\bm{P}=\left(
            \bm{\alpha}_1,\bm{\alpha}_2,\bm{\alpha}_3
        \right)$,由$\bm{AP}=\bm{PJ}$知\[
        \bm{A\alpha}_1=-\bm{\alpha}_1,\bm{A\alpha}_2=-\bm{\alpha}_2,\bm{A\alpha}_3=\bm{\alpha}_2-\bm{\alpha}_3
        \]于是$\bm{\alpha}_2$是$\bm{A}$关于特征值$-1$的特征向量,而$\bm{\alpha}_3$是$\left(
            \bm{A}+\bm{I}_3
        \right)\bm{x}=\bm{\alpha}_2$的一个解即广义特征向量.问题在于$\bm{\alpha}_2$的选取必然影响到$\bm{\alpha}_3$的选取.

        首先求得\begin{align*}
            \bm{A}+\bm{I}_3=\begin{pmatrix}
                3&6&-15\\1&2&-5\\1&2&-5
            \end{pmatrix}\longrightarrow\begin{pmatrix}
                1&2&-5\\0&0&0\\0&0&0
            \end{pmatrix}
        \end{align*}于是基础解系\[
        \bm{\beta}_1=\begin{pmatrix}
            -2\\1\\0
        \end{pmatrix},\bm{\beta}_2=\begin{pmatrix}
            5\\0\\1
        \end{pmatrix}
        \]试求解\[
        \left(
        \begin{array}{c:c}
            \bm{A}+\bm{I}_3&\bm{\beta}_2
        \end{array}
        \right)=\begin{pmatrix}
            3&6&-15&5\\1&2&-5&0\\1&2&-5&1
        \end{pmatrix}
        \]这是无解的.这并不是无解的问题,而是对$\bm{\alpha}_2$的选取不够好.取$\bm{\alpha}_2=k_1\bm{\beta}_1+k_2\bm{\beta}_2=\left(
            -2k_1+5k_2,k_1,k_2
        \right)'$并且此处应当有\[
        \mathrm{rank}\,\begin{pmatrix}
            3&6&-15&-2k_1+5k_2\\1&2&-5&k_1\\1&2&-5&k_2   
        \end{pmatrix}=1
        \]明显有$k_1=k_2$,于是不妨取$k_1=k_2=1$即$\bm{\alpha}_2=\begin{pmatrix}
            3\\1\\1
            \end{pmatrix}$,于是解出$\bm{\alpha}_3=\begin{pmatrix}
                1\\0\\0
            \end{pmatrix}$.而$\bm{\alpha}_1$取得与$\bm{\alpha}_2$线性无关即可,不妨取$\bm{\alpha}_1=\bm{\beta}_1=\begin{pmatrix}
                -2\\1\\0
            \end{pmatrix}$.于是所求\[
            \bm{P}= \begin{pmatrix}
                -2&3&1\\1&1&0\\0&1&0
            \end{pmatrix}
            \]
    \end{solution}
\end{brown}
\begin{red}
    \begin{remark}
        \textup{Jordan}标准型应用的第二个主题是如何求得带参数的矩阵的\textup{Jordan}标准型.主要有以下方法和思路：

        \begin{enumerate}[label={\textup{(\arabic*)}}]
            \item 选取特殊的子式求得行列式因子
            \item 计算度数以确定\textup{Jordan}块的个数
            \item 计算极小多项式以确定最大的\textup{Jordan}块的阶数
        \end{enumerate}
    \end{remark}
\end{red}
\begin{brown}
    \begin{example}
        设\[
        \bm{A}=\begin{pmatrix}
            a&1&1&\cdots&1\\
            &a&1&\cdots&1\\
            &&\ddots&\ddots&\vdots\\
            &&&\ddots&1\\
            &&&&a
        \end{pmatrix}
        \]
        求其\textup{Jordan}标准型.
    \end{example}
    \begin{solution}
        法一.
        \[
        \lambda\bm{I}-\bm{A}=\begin{pmatrix}
            \lambda-a&-1&-1&\cdots&-1\\
            &\lambda-a&-1&\cdots&-1\\
            &&\ddots&\ddots&\vdots\\
            &&&\ddots&-1\\
            &&&&\lambda-a  
        \end{pmatrix}
        \]注意到\[
        \left(
            \lambda\bm{I}-\bm{A}
        \right)\begin{pmatrix}
            1&2&\cdots&n-1\\
            1&2&\cdots&n-1
        \end{pmatrix}=\left(\lambda-a\right)^{n-1}
        \]设\begin{align*}
            \left(
                \lambda\bm{I}-\bm{A}
            \right)\begin{pmatrix}
                1&2&\cdots&n-1\\
                2&3&\cdots&n
            \end{pmatrix}=g\left(\lambda\right)
        \end{align*}并且$a$不是$g\left(\lambda\right)$的根即上述两子式互素.则根据\cref{def:行列式因子}有\[
        D_{n-1}\left(\lambda\right)=1,D_n\left(\lambda\right)=\left(\lambda-a\right)^n
        \]于是不变因子组\[
        1,\cdots,1,\left(\lambda-a\right)^n
        \]则\textup{Jordan}标准型为$\bm{J}_n\left(a\right).$

        法二.特征值$a$的代数重数为$n$,几何重数$=n-\mathrm{rank}\,\left(
            \bm{A}-a\bm{I}
        \right)=1$,于是\textup{Jordan}块只有一个$n$阶的即$\bm{J}_n\left(a\right)$.

        法三.$\bm{A}$的特征多项式为$\left(
            \lambda-a
        \right)^n$则极小多项式为$\left(
            \lambda-a
        \right)^k,k\leqslant n.$因为\[
        \bm{A}=a\bm{I}_n+\bm{N}+\bm{N}^2+\cdots+\bm{N}^{n-1}
        \]其中$\bm{N}=\bm{J}_n\left(0\right).$考虑$k=n-1$此时\begin{align*}
            \left(
                \bm{A}-a\bm{I}_n
            \right)^k&=\left(
                \bm{N}+\bm{N}^2+\cdots+\bm{N}^{n-1}
            \right)^{n-1}\\
            &=\bm{N}^{n-1}\neq \bm{O}
        \end{align*}于是极小多项式即为$\left(
            \lambda-a
            \right)^n$.则不变因子组为\[
            1,\cdots,1,\left(\lambda-a\right)^n\]于是\textup{Jordan}标准型为$\bm{J}_n\left(a\right)$.
    \end{solution}
\end{brown}
\begin{brown}
    \begin{example}
        设矩阵\[
        \bm{A}=\begin{pmatrix}
            1&&&\\
            a+2&1&&\\
            5&3&1&\\
            7&6&b+4&1
        \end{pmatrix}
        \]求其\textup{Jordan}标准型.
    \end{example}
    \begin{solution}
        明显特征值为$1$,代数重数$4$,几何重数\[4-\mathrm{rank}\,\left(
            \bm{A}-\bm{I}_4
        \right)=\begin{cases*}
            1&,$a\neq -2$且$b\neq -4\leadsto\bm{J}_4\left(1\right)$\\ 
            2&,$a=-2$或$b=-4\leadsto\mathrm{diag}\,\left\{
                \bm{J}_k\left(1\right),\bm{J}_l\left(1\right)
            \right\},1\leqslant k\leqslant l,k+l=4$
        \end{cases*}\]显然$2\leqslant l\leqslant 3.$
        
        于是需要估计$\bm{A}$的极小多项式（其可能是$\left(\lambda-1\right)^l,l=2,3$）.首先\begin{align*}
            \left(\bm{A}-\bm{I}_4\right)^2&=\begin{pmatrix}
                &&&\\
                &&&\\
                3\left(a+2\right)&&&\\
                6\left(a+2\right)+5\left(b+4\right)&3\left(b+4\right)&&
            \end{pmatrix}\\
            &=\begin{cases*}
                \bm{O}&,$a=-2$且$b=-4\leadsto\mathrm{diag}\,\left\{
                    \bm{J}_2\left(1\right),\bm{J}_2\left(1\right)
                \right\}$\\
                \neq \bm{O}&,$a\neq -2$或$b\neq -4\leadsto\mathrm{diag}\,\left\{
                    \bm{J}_1\left(1\right),\bm{J}_3\left(1\right)
                \right\}$
            \end{cases*}
        \end{align*}
    \end{solution}
\end{brown}
\begin{red}
\begin{remark}
    \textup{Jordan}标准型应用的第三个主题是关于循环子空间的应用.
\end{remark}
\end{red}
\begin{formal}
\begin{definition}[一般的循环子空间定义]\label{def:一般的循环子空间定义}
    设线性变换$\bm{\varphi}\in\mathcal{L}\left(
        V_{\mathbb{K}}^n
    \right),\bm{0}\neq\bm{\alpha}\in V$.由$\left\{
        \bm{\alpha},\bm{\varphi}\left(\bm{\alpha}\right),\bm{\varphi}^2\left(\bm{\alpha}\right),\cdots
    \right\}$这一族向量生成的子空间$C\left(
        \bm{\varphi},\bm{\alpha}
    \right)$称为线性变换$\bm{\varphi}$关于循环向量$\bm{\alpha}$的循环子空间.
\end{definition}
\end{formal}
\begin{green}
\begin{lemma}[循环子空间的基]\label{lemma:循环子空间的基}
    设$\dim\,C\left(\bm{\varphi},\bm{\alpha}\right)=m$则$\left\{
        \bm{\alpha},\bm{\varphi}\left(\bm{\alpha}\right),\cdots,\bm{\varphi}^{m-1}\left(\bm{\alpha}\right)
    \right\}$必定是$C\left(\bm{\varphi},\bm{\alpha}\right)$的一组基.
\end{lemma}\begin{Proof}
    设\[k=\max\,\left\{
        i\in\mathbb{Z}^+\mid \bm{\alpha},\bm{\varphi}\left(\bm{\alpha}\right),\cdots,\bm{\varphi}^{i-1}\left(\bm{\alpha}\right)\text{线性无关}
    \right\}\]即$\bm{\alpha},\bm{\varphi}\left(\bm{\alpha}\right),\bm{\varphi}^2\left(\bm{\alpha}\right),\cdots,\bm{\varphi}^{k-1}\left(\bm{\alpha}\right)$线性无关且$\bm{\alpha},\bm{\varphi}\left(\bm{\alpha}\right),\bm{\varphi}^2\left(\bm{\alpha}\right),\cdots,\bm{\varphi}^{k}\left(\bm{\alpha}\right)$线性相关.根据\cref{cor:加入向量的线性关系}知$\bm{\varphi}^k\left(\bm{\alpha}\right)$是$\bm{\alpha},\bm{\varphi}\left(\bm{\alpha}\right),\bm{\varphi}^2\left(\bm{\alpha}\right),\cdots,\bm{\varphi}^{k-1}\left(\bm{\alpha}\right)$的线性组合.

    下一步考虑数学归纳法证明$\forall i\geqslant k,\bm{\varphi}^i\left(\bm{\alpha}\right)$是$\bm{\alpha},\bm{\varphi}\left(\bm{\alpha}\right),\bm{\varphi}^2\left(\bm{\alpha}\right),\cdots,\bm{\varphi}^{k-1}\left(\bm{\alpha}\right)$的线性组合.$i=k$时已经证明,设$<i$时成立,下面考虑$=i$时,因为$\bm{\varphi}^{i}\left(\bm{\alpha}\right)=\bm{\varphi}^{i-1}\left(\bm{\varphi}\left(\bm{\alpha}\right)\right)$是$\bm{\alpha},\bm{\varphi}\left(\bm{\alpha}\right),\bm{\varphi}^2\left(\bm{\alpha}\right),\cdots,\bm{\varphi}^{k-1}\left(\bm{\alpha}\right)$的线性组合.于是证毕.

    于是$\bm{\alpha},\bm{\varphi}\left(\bm{\alpha}\right),\bm{\varphi}^2\left(\bm{\alpha}\right),\cdots,\bm{\varphi}^{k-1}\left(\bm{\alpha}\right)$是$C\left(\bm{\varphi},\bm{\alpha}\right)$的一组基并且维数为$k$即$m=\dim\,C\left(\bm{\varphi},\bm{\alpha}\right)=k$.
\end{Proof}
\end{green}
\begin{green}
\begin{corollary}[最小不变子空间]\label{cor:最小不变子空间}
    事实上,$C\left(
        \bm{\varphi},\bm{\alpha}
    \right)$是包含$\bm{\alpha}$的最小$\bm{\varphi}$-不变子空间.
\end{corollary}
\end{green}
\begin{green}
    \begin{corollary}[限制的表示矩阵]\label{cor:限制的表示矩阵}
设\[
\bm{\varphi}^m\left(\bm{\alpha}\right)=-a_0\bm{\alpha}-a_1\bm{\varphi}\left(\bm{\alpha}\right)-\cdots-a_{m-1}\bm{\varphi}^{m-1}\left(\bm{\alpha}\right)
\]令\[
g\left(x\right)=x^m+a_{m-1}x^{m-1}+\cdots+a_1x+a_0\in \mathbb{K}\left[x\right]
\]而限制$\bm{\varphi}\Big|_{C\left(
    \bm{\varphi},\bm{\alpha}
\right)}$在基$\left\{
    \bm{\alpha},\bm{\varphi}\left(\bm{\alpha}\right),\cdots,\bm{\varphi}^{m-1}\left(\bm{\alpha}\right)
\right\}$下的表示矩阵为\[
\bm{F}\left(g\left(\lambda\right)\right)=\begin{pmatrix}
    &&&&-a_0\\
    1&&&&-a_1\\
    &1&&&\vdots\\
    &&\ddots&&\vdots\\
    &&&1&-a_{m-1}
\end{pmatrix}
\]    
    \end{corollary}
\end{green}
\begin{brown}
    \begin{example}\,

        \begin{enumerate}[label={\textup{(\arabic*)}}]
            \item 考虑$\bm{\varphi}\left(\alpha\right)=\lambda_0\bm{\alpha},\lambda_0\in\mathbb{K},\bm{\alpha}\neq \bm{0}$的循环子空间$C\left(
                \bm{\varphi},\bm{\alpha}
            \right)=L\left(
                \bm{\alpha}
            \right)$.
            
            \item 进一步地,考虑$\bm{\varphi}\left(\bm{\alpha}_i\right)=\lambda_i\bm{\alpha}_i,\lambda_1,\lambda_2,\cdots,\lambda_n\in\mathbb{K}$且互异,$\bm{\alpha}_i\neq\bm{0}$.令$\bm{\alpha}=\bm{\alpha}_1+\bm{\alpha}_2+\cdots+\bm{\alpha}_n$,容易证明$C\left(
                \bm{\varphi},\bm{\alpha}
            \right)=V$.
        \end{enumerate}
    \end{example}
\end{brown}
\begin{brown}
    \begin{example}\label{ex:全空间分解}
        设$\bm{\varphi}\in\mathcal{L}\left(
            V_{\mathbb{K}}^n
        \right)$的不变因子为\[
        1,\cdots,1,d_1\left(\lambda\right),\cdots,d_k\left(\lambda\right)
        \]根据有理标准型理论知存在$V$的一组基$\left\{
            \bm{e}_1,\bm{e}_2,\cdots,\bm{e}_n
        \right\}$使得$\bm{\varphi}$在这组基下的表示矩阵为\[
        \bm{F}= \mathrm{diag}\,\left\{
            \bm{F}\left(d_1\left(\lambda\right)\right),\bm{F}\left(d_2\left(\lambda\right)\right),\cdots,\bm{F}\left(d_k\left(\lambda\right)\right)
        \right\}
        \]因为表示矩阵为多项式友阵说明是循环子空间,于是全空间一定可以分解为循环子空间的直和\[
    V=C\left(
        \bm{\varphi},\bm{e}_{i_1}
     \right)\oplus C\left(
        \bm{\varphi},\bm{e}_{i_2}
        \right)\oplus\cdots\oplus C\left(
            \bm{\varphi},\bm{e}_{i_k}
        \right)
        \]
    \end{example}
\end{brown}
\begin{brown}
    \begin{example}\label{ex:全空间分解2}
        设$\bm{\varphi}\in\mathcal{L}\left(
            V_{\mathbb{C}}^n
        \right)$的初等因子为\[
        \left(\lambda-\lambda_1\right)^{r_1},\left(\lambda-\lambda_2\right)^{r_2},\cdots,\left(\lambda-\lambda_k\right)^{r_k}
        \]并且存在一组基$\left\{
            \bm{e}_1,\bm{e}_2,\cdots,\bm{e}_n
        \right\}$使得$\bm{\varphi}$在这组基下的表示矩阵为\[
        \bm{J}= \mathrm{diag}\,\left\{
            \bm{J}_{r_1}\left(\lambda_1\right),\bm{J}_{r_2}\left(\lambda_2\right),\cdots,\bm{J}_{r_k}\left(\lambda_k\right)
        \right\}
        \]根据\cref{thm:全空间关于Jordan标准型的两种直和分解}知\[
        V=C\left(
            \bm{\varphi}-\lambda_1\bm{I}_V,\bm{e}_{i_1}
        \right)\oplus C\left(
            \bm{\varphi}-\lambda_2\bm{I}_V,\bm{e}_{i_2}
        \right)\oplus\cdots\oplus C\left(
            \bm{\varphi}-\lambda_k\bm{I}_V,\bm{e}_{i_k}
        \right)
        \]
    \end{example}
\end{brown}
\begin{red}
    \begin{remark}
        \cref{ex:全空间分解}和\cref{ex:全空间分解2}都是关于全空间的分解,但是
        \cref{ex:全空间分解}是关于同一个线性变换而\cref{ex:全空间分解2}不是.
    \end{remark}
\end{red}
\begin{formal}
    \begin{proposition}[唯一确定可交换的线性变换]\label{prop:唯一确定可交换的线性变换}
        设$n$维循环空间$V=C\left(
            \bm{\varphi},\bm{\alpha}
        \right),\bm{\psi}\in\mathcal{L}\left(
            V
        \right)$且$
        \bm{\varphi\psi}=\bm{\psi\varphi}
        $则\begin{enumerate}[label={\textup{(\arabic*)}}]
            \item $\bm{\psi}$由$\bm{\psi}\left(\bm{\alpha}\right)$唯一确定
            \item $\exists g\left(x\right)\in\mathbb{K}\left[x\right]\,\textup{s.t.}\bm{\psi}=g\left(\bm{\varphi}\right)$
        \end{enumerate}
    \end{proposition}
    \begin{Proof}$(1)$
        因为$V=C\left(
            \bm{\varphi},\bm{\alpha}
        \right)$的基为$\left\{
            \bm{\alpha},\bm{\varphi}\left(\bm{\alpha}\right),\cdots,\bm{\varphi}^{n-1}\left(\bm{\alpha}\right)
        \right\}$,$\bm{\psi}$由其在基上的作用唯一决定而$\bm{\psi}\left(
            \bm{\varphi}^i\left(
                \bm{\alpha}
            \right)
        \right)=\bm{\varphi}^i\left(
            \bm{\psi}\left(
                \bm{\alpha}
            \right)
        \right),\forall 0\leqslant i\leqslant n-1$.于是$\bm{\psi}$由$\bm{\psi}\left(
            \bm{\alpha}
        \right)$唯一确定.

        $(2)$因为\[
        \bm{\psi}\left(\bm{\alpha}\right)=
        a_0\bm{\alpha}+a_1\bm{\varphi}\left(\bm{\alpha}\right)+\cdots+a_{n-1}\bm{\varphi}^{n-1}\left(\bm{\alpha}\right)
        \]令\[
        g\left(x\right)=a_0+a_1x+\cdots+a_{n-1}x^{n-1}\in\mathbb{K}\left[x\right]
        \]于是$\bm{\psi}$和$g\left(\bm{\varphi}\right)$是两个均可与$\bm{\varphi}$交换的线性变换,且$(1)$表明它们由$\bm{\alpha}$处的值唯一确定,于是由$\bm{\psi}\left(
            \bm{\alpha}
        \right)=g\left(
            \bm{\varphi}
        \right)\left(
            \bm{\alpha}
        \right)\Longrightarrow \bm{\psi}=g\left(
            \bm{\varphi}
        \right).$
    \end{Proof}
\end{formal}
\begin{green}
    \begin{corollary}[相关推论]\label{cor:相关推论}
        \,
        
        \begin{enumerate}[label={\textup{(\arabic*)}}]
            \item 若$\bm{A}=\bm{F}\left(g\left(\lambda\right)\right),\bm{AB}=\bm{BA}$,则$\bm{B}=f\left(\bm{A}\right)$
            \item 若$\bm{A}=\bm{J}_n\left(\lambda_0\right),\bm{AB}=\bm{BA}$,则$\bm{B}=g\left(\bm{A}\right)$
        \end{enumerate}
    \end{corollary}
\end{green}
\newpage
\section{矩阵函数}
\subsection{级数相关知识}
\begin{red}
    \begin{remark}
        定义矩阵函数至少有三种方法：其一,将函数作幂级数展开并利用\textup{Jordan}标准型理论定义；其二,利用多项式定义；其三,利用\textup{Cauchy}积分定理定义.这里采用第一种定义.
    \end{remark}
\end{red}
\begin{formal}
    \begin{definition}[级数收敛的定义]\label{def:级数收敛的定义}
        无穷级数$\displaystyle \sum_{n=1}^{\infty}a_n$收敛当且仅当部分和数列$\left\{s_n\right\}$收敛,其中$s_n=a_1+a_2+\cdots+a_n.$若\[
        \lim_{n\to \infty}s_n=S
        \]则记级数$\displaystyle \sum_{n=1}^{\infty}a_n=S$.
    \end{definition}
\end{formal}
\begin{formal}
    \begin{theorem}[Cauchy审敛法]\label{thm:Cauchy审敛法}
        考虑正项级数$\displaystyle \sum_{n=1}^{\infty}a_n$,令$\displaystyle r=\overline{\lim_{n\to\infty}}\sqrt[n]{a_n}$则\begin{enumerate}[label={\textup{(\arabic*)}}]
            \item $\displaystyle \sum_{n=1}^{\infty}a_n$收敛当且仅当$r<1$
            \item $\displaystyle \sum_{n=1}^{\infty}a_n$发散当且仅当$r>1$
            \item 当$r=1$时无法判断
        \end{enumerate}
    \end{theorem}
\end{formal}
\begin{formal}
    \begin{definition}[幂级数的收敛半径的定义]\label{def:幂级数的收敛半径的定义}
        考虑幂级数$\displaystyle \sum_{n=0}^{\infty}a_nx^n$,令$\displaystyle A=\overline{\lim_{n\to\infty}}\sqrt[n]{\left|a_n\right|}$,定义该幂级数收敛半径\[
        R=\begin{cases*}
            +\infty&,$A=0$\\
            \dfrac{1}{A}&,$0<A<+\infty$\\
            0&,$A=+\infty$
        \end{cases*}
        \]
    \end{definition}
\end{formal}
\begin{formal}
    \begin{theorem}[Cauchy-Hadamard定理]\label{thm:Cauchy-Hadamard定理}
        在\cref{def:幂级数的收敛半径的定义}的基础上,幂级数$\displaystyle \sum_{n=0}^{\infty}a_nx^n$的敛散性为\begin{enumerate}[label={\textup{(\arabic*)}}]
            \item $\left|x\right|<R$时,$\displaystyle \sum_{n=0}^{\infty}a_nx^n$绝对收敛
            \item $\left|x\right|>R$时,$\displaystyle \sum_{n=0}^{\infty}a_nx^n$发散
            \item $\left|x\right|=R$时,敛散性不确定
        \end{enumerate}
    \end{theorem}
\end{formal}
\begin{formal}
    \begin{theorem}[逐项可导性]\label{thm:逐项可导性}
        幂级数$\displaystyle f\left(x\right)=\sum_{n=0}^{\infty}a_nx^n$在$\left|x\right|<R$时逐项可导即\[
        f'\left(x\right)=\frac{\mathrm{d}}{\mathrm{d}x}\left(
            \sum_{n=0}^{\infty}a_nx^n
        \right)=\sum_{n=0}^{\infty}\frac{
            \mathrm{d}
        }{\mathrm{d}x}\left(
            a_nx^n
        \right)=\sum_{n=0}^{\infty}na_nx^{n-1}
        \]并且$\displaystyle f'\left(x\right)=
        \sum_{n=0}^{\infty}na_nx^{n-1}
         $的收敛半径也是$R$.
    \end{theorem}
\end{formal}
\begin{red}
\begin{remark}
    下面考虑复幂级数
    \[
    f\left(z\right)=\sum_{n=0}^{\infty}a_nz^n
    ,a_i\in\mathbb{C},z\in\mathbb{C}
    \]值得注意的是$\left|z\right|$表示复数的模长而不是实数的绝对值.其余表述与实幂级数类似.
\end{remark}
\end{red}
\begin{formal}
    \begin{proposition}[常见复幂级数及其收敛半径]\label{prop:常见复幂级数及其收敛半径}
        \,

        \begin{enumerate}[label={\textup{(\arabic*)}}]
            \item $\displaystyle \mathrm{e}^z=1+z+\frac{z^2}{2!}+\frac{z^3}{3!}+\cdots,R=+\infty$
            \item $\displaystyle \sin z=z-\frac{z^3}{3!}
            +\frac{z^5}{5!}-\frac{z^7}{7!}+\cdots,R=+\infty
            $
            \item $\displaystyle \cos z=1-\frac{z^2}{2!}+\frac{z^4}{4!}-\frac{z^6}{6!}+\cdots,R=+\infty$
            \item $\ln \left(1+z\right)=
            z-\frac{z^2}{2}+\frac{z^3}{3}-\frac{z^4}{4}+\cdots,R=1
            $
        \end{enumerate}
    \end{proposition}
\end{formal}
\subsection{矩阵函数}
考虑$f\left(z\right)=a_0+a_1z+\cdots+a_pz^p\in \mathbb{C}\left[z\right]$,矩阵$\bm{A}\in M_n\left(\mathbb{C}\right)$则\[
f\left(\bm{A}\right):=a_0\bm{I}_n+a_1\bm{A}+\cdots+a_p\bm{A}^p
\]

设非异阵$\bm{P}$使得\[
\bm{P}^{-1}\bm{A}\bm{P}=\bm{J}= \mathrm{diag}\,\left\{
    \bm{J}_{r_1}\left(\lambda_1\right),\bm{J}_{r_2}\left(\lambda_2\right),\cdots,\bm{J}_{r_k}\left(\lambda_k\right)
\right\}
\]并且$f\left(\bm{A}\right)=\bm{P}f\left(\bm{J}\right)\bm{P}^{-1}=\bm{P}\mathrm{diag}\,\left\{
    f\left(
      \bm{J}_{r_1}\left(\lambda_1\right)
    \right),f\left(
        \bm{J}_{r_2}\left(\lambda_2\right)
    \right),\cdots,f\left(
        \bm{J}_{r_k}\left(\lambda_k\right)
    \right)
\right\}\bm{P}^{-1}$.

于是考虑每一个Jordan块,由\cref{thm:Jordan-Chevalley分解定理}知有分解$\bm{J}_{r_i}\left(\lambda_i\right)=\lambda_i\bm{I}+\bm{N}$,其中$\bm{N}=\bm{J}_{r_i}\left(0\right),\bm{N}^{r_i}=\bm{O}$ .

考虑$f\left(z\right)$在$\lambda_i$处的Taylor展开\[
f\left(z\right)=f\left(\lambda_i\right)+\frac{f'\left(\lambda_i\right)}{1!}\left(z-\lambda_i\right)+\frac{f''\left(\lambda_i\right)}{2!}\left(z-\lambda_i\right)^2+\cdots+\frac{
    f^{\left(p\right)}\left(\lambda_i\right)
}{p!}\left(z-\lambda_i\right)^p
\]代入$\bm{J}_{r_i}\left(\lambda_i\right)=\lambda_i\bm{I}+\bm{N}$得到\begin{align*}
    f\left(
        \bm{J}_{r_i}\left(
        \lambda_i
        \right)
    \right)&=f\left(\lambda_i\right)\bm{I}+\frac{f'\left(\lambda_i\right)}{1!}\bm{N}+\frac{f''\left(\lambda_i\right)}{2!}\bm{N}^2+\cdots+\frac{
        f^{\left(r_i-1\right)}\left(\lambda_i\right)
    }{\left(r_i-1\right)!}\bm{N}^{r_i-1}\\
    &=\begin{pmatrix}
        f\left(\lambda_i\right)&\cfrac{f'\left(\lambda_i\right)}{1!}&\cfrac{f''\left(\lambda_i\right)}{2!}&\cdots&\cfrac{f^{\left(r_i-1\right)}\left(\lambda_i\right)}{\left(r_i-1\right)!}\\
        &f\left(\lambda_i\right)&\cfrac{f'\left(\lambda_i\right)}{1!}&\cdots&\cfrac{f^{\left(r_i-2\right)}\left(\lambda_i\right)}{\left(r_i-2\right)!}\\
        &&f\left(\lambda_i\right)&\cdots&\cfrac{f^{\left(r_i-3\right)}\left(\lambda_i\right)}{\left(r_i-3\right)!}\\
        &&&\ddots&\vdots\\
        &&&&f\left(\lambda_i\right)
    \end{pmatrix}
\end{align*}
下面给出矩阵幂级数的收敛\begin{formal}
    \begin{definition}[矩阵幂级数收敛的定义]\label{def:矩阵幂级数收敛的定义}
        设矩阵序列$\left\{\bm{A}_p\right\}$,其中\[
        \bm{A}_p=\begin{pmatrix}
            a_{11}^{\left(p\right)}&a_{12}^{\left(p\right)}&\cdots&a_{1n}^{\left(p\right)}\\
            a_{21}^{\left(p\right)}&a_{22}^{\left(p\right)}&\cdots&a_{2n}^{\left(p\right)}\\
            \vdots&\vdots&&\vdots\\
            a_{n1}^{\left(p\right)}&a_{n2}^{\left(p\right)}&\cdots&a_{nn}^{\left(p\right)}
        \end{pmatrix}
        \]和矩阵$\bm{B}=\left(b_{ij}\right)_{n\times n}.$若$\forall 1\leqslant i,j\leqslant n,$数列$\left\{
            a_{ij}^{\left(p\right)}
        \right\}$收敛到$b_{ij}$即$\displaystyle \lim_{p\to\infty}a_{ij}^{\left(p\right)}=b_{ij},$则称矩阵序列$\left\{\bm{A}_p\right\}$收敛到矩阵$\bm{B}$,记作$\displaystyle \lim_{p\to\infty}\bm{A}_p=\bm{B},$否则称矩阵序列$\left\{\bm{A}_p\right\}$发散.
    \end{definition}
\end{formal}
\begin{brown}
    \begin{example}
        \[
        \bm{A}_p=\begin{pmatrix}
            1+\cfrac{1}{p}&\cfrac{1}{p^2}\\
            2&3
        \end{pmatrix}
        \]收敛于\[
        \bm{B}=\begin{pmatrix}
            1&0\\
            2&3
        \end{pmatrix}
        \]
    \end{example}
\end{brown}
\begin{formal}
    \begin{definition}[部分和定义收敛]\label{def:部分和定义收敛}
        设复幂级数$\displaystyle f\left(z\right)=\sum_{i=0}^{\infty}a_iz^i,$其部分和\[
        f_p\left(z\right)=a_0+a_1z+\cdots+a_pz^p
        \]设$\bm{A}\in M_n\left(\mathbb{C}\right)$得到矩阵序列$\displaystyle \left\{
            f_p\left(\bm{A}\right)
        \right\}_{p=0}^{\infty}$.若该矩阵序列收敛到矩阵$\bm{B}$,则矩阵幂级数\[
        \sum_{i=0}^{\infty}a_i\bm{A}^i=a_0\bm{I}_n+a_1\bm{A}+a_2\bm{A}^2+\cdots
        \]收敛到$\bm{B}$,记作$f\left(\bm{A}\right)=\bm{B}.$
    \end{definition}
\end{formal}
\begin{formal}
    \begin{definition}[矩阵函数的定义]\label{def:矩阵函数的定义}
        将\cref{def:部分和定义收敛}中取未定元$\bm{X}$则得到一般的矩阵函数\[
        f\left(\bm{X}\right)=\sum_{i=0}^{\infty}a_i\bm{X}^i=a_0\bm{I}_n+a_1\bm{X}+a_2\bm{X}^2+\cdots
        \]
    \end{definition}
\end{formal}
\begin{formal}
    \begin{theorem}[矩阵幂级数的收敛判定]\label{thm:矩阵幂级数的收敛判定}
        设复幂级数$\displaystyle f\left(z\right)=\sum_{i=0}^{\infty}a_iz^i$,收敛半径$R$,则\begin{enumerate}[label={\textup{(\arabic*)}}]
            \item $f\left(\bm{X}\right)$收敛当且仅当对于任意非异阵$\bm{P}$,$f\left(
                \bm{P}^{-1}\bm{X}\bm{P}
            \right)$都收敛
            \item 设\[\bm{X}=\mathrm{diag}\,\left\{
                \bm{X}_1,\bm{X}_2,\cdots,\bm{X}_k
            \right\}\],则$f\left(\bm{X}\right)$收敛当且仅当$f\left(
                \bm{X}_1
            \right),f\left(
                \bm{X}_2
            \right),\cdots,f\left(
                \bm{X}_k
            \right)$收敛
            \item 对\textup{Jordan}块$\bm{J}_{r_i}\left(\lambda_i\right)$,若$\left|\lambda_i\right|<R$,则$f\left(
                \bm{J}_{r_i}\left(\lambda_i\right)
            \right)$收敛到\[
            f\left(
                \bm{J}_{r_i}\left(\lambda_i\right)
            \right)=\begin{pmatrix}
                f\left(\lambda_i\right)&\cfrac{f'\left(\lambda_i\right)}{1!}&\cfrac{f''\left(\lambda_i\right)}{2!}&\cdots&\cfrac{f^{\left(r_i-1\right)}\left(\lambda_i\right)}{\left(r_i-1\right)!}\\
                &f\left(\lambda_i\right)&\cfrac{f'\left(\lambda_i\right)}{1!}&\cdots&\cfrac{f^{\left(r_i-2\right)}\left(\lambda_i\right)}{\left(r_i-2\right)!}\\
                &&f\left(\lambda_i\right)&\cdots&\cfrac{f^{\left(r_i-3\right)}\left(\lambda_i\right)}{\left(r_i-3\right)!}\\
                &&&\ddots&\vdots\\
                &&&&f\left(\lambda_i\right)
            \end{pmatrix}
            \]
        \end{enumerate}
    \end{theorem}
    \begin{Proof}
        
    \end{Proof}
\end{formal}